We recall the notion of Frobenius extensions, Demazure operators, and Frobenius hypercubes. These will be used in the next chapter to construct a diagrammatic category as in \cite{EWSFrob}. 
This chapter ends with various calculations involving Demazure operators for the deformed affine realization. These calculations will be used in section \ref{sec well-defined-ness of the functor}.
%==========================================================================
\subsection{Cubes of Frobenius Extensions} \label{subsec- cubes of frob}
%==========================================================================
 % In this section, all rings will be assumed to be commutative, and all gradings will be over $\Z$.

% \begin{defn} A \emph{(commutative) Frobenius extension} is an extension of commutative rings $\iota :  A \hookrightarrow B$, where $B$ is finitely generated free $A$-module, equipped with an
% $A$-linear map $\dd = \dd_A^B : B \to A$, called the \emph{(Frobenius) trace}. The trace is required to be \emph{non-degenerate}, meaning the following map is an isomorphism:
% $$B \rightarrow \Hom_A(B,A): b \mapsto (b'\mapsto \dd(bb')).$$

\begin{defn} A \emph{graded (commutative) Frobenius extension of degree d} is an extension of commutative graded rings $\iota :  A \hookrightarrow B$, where $B$ is a finitely generated free graded $A$-module, equipped with an
$A$-linear map $\dd = \dd_A^B : B \to A$ of degree $-2d$, called the \emph{(Frobenius) trace}. In other words, $\dd$ is a homomorphism of graded $A$-modules $\dd: B \rightarrow A(-2d)$. The trace is required to be \emph{non-degenerate}, meaning the following map is an isomorphism:
$$B \rightarrow \Hom_A(B,A): b \mapsto (b'\mapsto \dd(bb')).$$
% as before after forgetting the gradings.  
% We shall use the notation $\deg(\dd^B_A)$ to refer to the degree of the graded Frobenius extension $A \hookrightarrow B$. \BE{not my favorite notation, since that would seem like a negative number to me. Once you define $\Delta$ you could use $\deg(\Delta)$...}\PT{I think it might be better to just remove this sentence completely since I  end up not using this notation much.}

\end{defn}
\begin{remark}\label{rmk-non-degeneracy and dual basis}
The non-degeneracy of $\dd$ is equivalent
to requiring that $B$ admits a \emph{pair of dual bases}. This is a pair of bases
$\{x_i\}$ and $\{y_i\}$ for $B$ as an $A$-module, such that $\dd(x_i y_j)=\delta_{i j}$. The equivalence between these two notions is familiar from linear algebra. 
Within $\Hom_A(B,A)$ there is a dual basis $\{x_i^*\}$ to the basis $\{x_i\}$, and one takes $y_{i}$ to be the preimage of $x_{i}^*$ under $B \to \Hom_A(B,A)$. 
If $\dd$ is non-degenerate then any basis $\{x_{i}\}$ admits a unique dual basis.

If we start with a homogeneous basis $\{x_i\}\subset B$, then the dual basis $\{y_i\}$ is also homogeneous. After all, if $\{y_{i}\}$ is the dual basis to $\{x_{i}\}$, then the appropriate homogeneous components of $\{y_{i}\}$ are also dual to $\{x_{i}\}$, but the dual basis is unique.


    
    % \begin{itemize}
    %     \item  If $\dd$ is non-degenerate, then for any basis $\{x_{\al}\}$ of $B$ as an $A$-module (which exists since $B$ is assumed to be a free $A$-module), we have the dual basis $\{x^*_{\al}\} \subset \Hom_A(B,A)$. 
    %     Taking $y_{\al}$ to be the preimage of $x^*_{\al}$ via the isomorphism $B \rightarrow \Hom_A(B,A)$ induced by $\dd$, we have:
    %     $$\dd(y_{\al}x_{\beta})=x^*_{\al}(x_{\beta})=\delta_{\al \beta}.$$
    %     \item If $\{x_{\al}\}, \{y_{\al}\}$ form a pair of dual basis, then we see that the map $$B \rightarrow \Hom_A(B,A): b \mapsto (b'\mapsto \dd(bb'))$$ takes $y_\al$ to $x_\al^*$ where $\{x_\al^*\} \subset \Hom_A(B,A)$ is the dual basis to $\{x_\al\}$. 
    %     In other words, the map $B \rightarrow \Hom_A(B,A)$ takes an $A$-basis of $B$ to an $A$-basis  of $\Hom_A(B,A)$ and is hence an isomorphism.
    % \end{itemize}


% \PT{I think this is a nice exercise for the reader- if $\deg(x_\al)=r_\al$, then the component $y_\al':=(y_\al)_{d-r_\al}$ also satisfies $\dd(x_\al y_\beta')=\delta_{\al \beta}$. } \BE{Ok, I like this exercise, but the rest of the remark is relatively obvious. I might just retool this remark so it talks about the graded case, and say this exercise aloud: that if $x$ is a homogeneous basis and $y$ is any dual basis then the appropriate homogeneous componetns of $y$ form a homogeneous dual basis.}

    % It follows that there exists \emph{dual basis} $\{x_\alpha\}$ and $\{y_\alpha\}$ in $B$, such that $\dd(x_\alpha y_\beta)=\delta_{\alpha \beta}$, which can be assumed to be \emph{homogeneous}, i.e., each of the $x_\al, y_\al$, can be assumed to be live purely in one of the graded components.
\end{remark}

The data of a Frobenius extension also equips one with a comultiplication map $\Delta^B_A : B \to B \otimes_A B$, which is a $B$-bimodule map.  It sends $1 \mapsto \sum_\alpha x_i \otimes
y_i$, an element which is independent of the choice of dual bases. If $A \hookrightarrow B$ is a graded Frobenius extension of degree $d$, then the comultiplication map $\Delta^B_A$ has degree $2d$.  

\begin{defn}\label{defn-product coproduct element}
    The \emph{product-coproduct} element $\mu^B_A$ of a Frobenius extension $A \hookrightarrow B$ is defined to be 
    $$\mu_A^B:=m^B_A\circ\Delta^B_A (1)$$
    where $m^B_A:B\otimes_AB \rightarrow B$ denotes the multiplication map of $B$.
\end{defn}

\begin{lemma}
The element $\partial^B_A(\mu^B_A) \in A$ is a scalar, equal to the rank (not the graded rank) of $B$ as an $A$-module. \end{lemma}

\begin{proof} This follows directly from the definition of dual bases. \end{proof}


Being a Frobenius extension is a structure (the choice of $\dd$), not just a property. However, the choice of $\dd$ is unique up to a choice of a unit in $B$. 

\begin{lemma} \label{lem Frobenius structures unique up to unit}
Suppose $\dd, \dd': B \rightarrow A$ are two Frobenius traces. Then there exists a unit $b_0 \in B$ such that $\dd'(b) = \dd(b_0 b)$ for all $b \in B$. If $\mu$ and $\mu'$ represent $\mu^B_A$ for these two choices of Frobenius structure, then $\mu = b_0 \mu'$.
    
Thus the product-coproduct element is independent of the choice of Frobenius trace up to unit, and if $\mu = \mu'$ then $\dd = \dd'$. \end{lemma} 

\begin{proof} Via the isomorphism 
    $$B \rightarrow \Hom_A(B,A): b \mapsto (b' \mapsto \dd(bb')),$$ we have that $\dd' \in \Hom_A(B,A)$ corresponds to some $b_0\in B$, so that $\dd'(b)=\dd(b_0b)$ for any $b \in B$. We leave the reader to verify that $b_0$ is a unit, and the remaining statements.
\end{proof}

    
    % Similarly, let $b_0'\in B$ correspond to $\dd$ via the isomorphism $B\cong\Hom_A(B,A)$ defined by $\dd'$, so that $\dd(b)=\dd'(b_0'b)$.
    % We then have that $\dd(b)=\dd'(b_0'b)=\dd(b_0b_0'b)$ for any $b \in B$, so that $b_0b_0'=1$ (since  via the isomorphism $B \cong \Hom_A(B,A)$ defined by $\dd$, 1 and $b_0b_0'$ both correspond to $\dd\in \Hom_A(B,A)$).
    % Hence we have that $b_0, b_0'$ are units in $B$, and that $\dd, \dd'$ are the same up to these units in $B$.


Of note is the case when $A \hookrightarrow B$ induces an equality $A^\times = B^\times$, because then $\dd' = a_0 \dd$ for $a_0 \in A^\times$. We say $\dd'$ is a rescaling of $\dd$. For example, in this paper it happens frequently that $A$ and $B$ are $\kk$-algebras and $A^\times=B^\times=\kk^\times$.

Note that if $A \subset B$ and $B \subset C$ are Frobenius
extensions, then $A \subset C$ is a Frobenius extension as well, with trace $\dd^C_A = \dd^B_A \dd^C_B$. It is proven in \cite[(2.5)]{EWSFrob} that \begin{equation} \label{eq:mumult} \mu^C_A = \mu^C_B \mu^B_A. \end{equation}


\begin{defn} A \emph{hypercube of (graded) Frobenius extensions} or a \emph{Frobenius hypercube} will be the following datum. \begin{itemize} \item
A finite set $\Gamma$. We also use $\Gamma$ to designate the entire datum. We consider the hypercube with vertices labelled by subsets of
$\Gamma$. An edge in this hypercube corresponds to $I \setminus \gamma \subset I$ for some $\gamma \in I \subset \Gamma$, and parallel
edges correspond to the same $\gamma$. \item A (contravariant) assignment of graded rings $R^I$ to vertices in the cube, so that $I \subset J
\implies R^J \subset R^I$. \item For each edge, a trace map $\dd^{I \setminus \gamma}_I : R^{I \setminus \gamma} \to R^I$ making $R^I \hookrightarrow R^{I
\setminus \gamma}$ into a graded Frobenius extension. 

\end{itemize}
We call this hypercube \emph{compatible} if, for every square $I \subset
J,J' \subset K$ in the hypercube, we have $\dd^J_K \dd^I_J = \dd^{J'}_K \dd^I_{J'}$. In this case,
there is a well-defined map $\dd^I_J : R^I \to R^J$ for every $I \subset J$, which endows the extension $R^J \subset R^I$ with the
structure of a Frobenius extension. We assume henceforth that every hypercube of Frobenius extensions is compatible.

We usually refer to $R^{\emptyset}$ as just $R$. Similarly we write $\dd_I := \dd^\emptyset_I$.

% Similarly, one defines a \emph{graded Frobenius hypercube} by replacing rings with graded rings and Frobenius extensions with graded Frobenius extensions.

A \emph{partial hypercube of Frobenius extensions} or a \emph{partial Frobenius hypercube} is a finite set $\Gamma$ and a non-empty collection $\mathcal{F}$ of subsets of $\Gamma$ (i.e. vertices in the hypercube), together with the data of graded rings and Frobenius extensions as above but only for vertices within the subset $\mathcal{F}$. We require that $J \in \mathcal{F}$ and $I \subset J$ implies $I \in \mathcal{F}$.  Note that $\emptyset$ is always in $\mc{F}$. For the purposes of this paper, we shall use the term Frobenius hypercube to also to refer to partial Frobenius hypercubes.
\end{defn}

For a Frobenius hypercube and $I \subset J \in \mathcal{F}$, let $\mu^I_J \in R^I$ denote the product-coproduct element for the Frobenius extension $R^J \subset R^I$, and $\mu_I = \mu^{\emptyset}_I$. By \eqref{eq:mumult} we have
\begin{equation} \label{eq:mumult2} \mu^I_J = \frac{\mu_J}{\mu_I}.\end{equation}
It is straightforward to see that \eqref{eq:mumult2} is equivalent to compatibility.

For sake of simplicity, assume that all rings in the hypercube are $\Bbbk$-algebras and that their unit groups agree with $\Bbbk^\times$.  By \cref{lem Frobenius structures unique up to unit}, each individual Frobenius extension $I \subset J$ in the cube is uniquely determined up to a scalar, and this scalar is visible in the choice of $\mu^I_J$. When defining a new Frobenius hypercube structure by rescaling the various elements $\mu^I_J$, constraints are placed on these scaling factors by \eqref{eq:mumult2}. Ultimately, one can freely rescale the elements $\mu_I \in R$ (resp. the trace maps $\partial_I$) for each $I \in \mathcal{F}$, and there will be a unique way to extend this rescaling to the entire hypercube so that compatibility will be maintained.

Fix a Coxeter system $(W,S)$ with a realization, and let $R$ be as in \cref{subsec- SSBim}. 
Let $\Gamma = S$, and $\mathcal{F}$ be the collection of finitary subsets of $S$. Under some assumptions on the realization, there will be a Frobenius hypercube with rings $R^I$ associated to each $I \in \mathcal{F}$. For example, when $W$ is a Weyl group with its geometric realization over $\mathbb{R}$, there is a standard Frobenius hypercube structure where $\mu_I$ is the product of the positive roots for the parabolic subgroup $W_I$. However, as we have noted in \S\ref{subsec-balanced}, some realizations do not have a well-behaved notion of positive roots. Consequently, the choice of product-coproduct elements $\mu_I$, or equivalently the choice of trace maps $\partial_I$, will be subtle. To construct the trace maps, we need to study Demazure operators.
 
% We shall discuss a large class of examples in section \ref{subsec-frob realizations} which arise from ``nice" realizations of Coxeter groups  (see \cref{example nice realizations are frobenius}).
% The motivating examples \cite[Example 1.4, Example 1.5]{EWSFrob} will be special cases of these. 



%==========================================================================
\subsection{Demazure operators}
\label{subsec-demazures}
%==========================================================================

For the rest of this section we fix a Coxeter system $(W,S)$, with a realization $(\Lambda,  \Lambda^\vee,\{\alpha_s\},\{\alpha_s^\vee\})$ over an integral domain $\kk$. We have the graded $\kk$-algebra $R:= \Sym_{\kk} \Lambda$, where $\Lambda$ is in degree 2, with its natural $W$-action.

We define \emph{the Demazure operator} or \emph{the divided difference operator} $\dd_s$ associated to a simple reflection $s\in S$ to be a $\kk$-linear operator on $R$:
$$\dd_s(f)= \frac{f-s(f)}{\al_s} \text{\hspace{20 pt}for $f\in R$}.$$
Note that $\dd_s$ may also be defined as the unique $\kk-$linear operator on $R$ satisfying the following properties:
\begin{itemize}
    \item $\dd_s|_{\Lambda}$ coincides with the functional defined by $\al_s^\vee$, i.e., $\dd_s(\lambda)=\langle\al_s^\vee,\lambda\rangle$ for $\lambda \in \Lambda$.
    \item $\dd_s$ satisfies the twisted Leibniz rule:
    \begin{equation}\label{eq twisted leibniz}
        \dd_s(fg)=\dd_s(f)g+s(f)\dd_s(g).
    \end{equation}
\end{itemize}

\begin{lemma}\label{lemma demazure operators coxeter relations}
    (\cite[Claim A.7]{ECathedral}) The Demazure operators $\dd_s, \dd_t$ associated to simple reflections $s,t \in S$ satisfy the following relations.
    \begin{enumerate}
        \item $\dd_s\dd_s=0$.
        % \item If $st=ts$, then $\dd_s\dd_t=\dd_t\dd_s$.
        \item For $m_{st} < \infty$, 
    \begin{equation} \label{eq:rescaled braid} \lambda \dd_s \dd_t \dd_s \cdots = \dd_t \dd_s \dd_t \cdots \end{equation} (both expressions alternate and have length $m_{st}$) for some unit $\lambda \in \kk^\times$. Moreover, $\lambda=1$ when $m_{st}$ is even.
    \end{enumerate}
The last property states that the braid relations hold up to scalar. This unit $\lambda$ agrees with the unit $\lambda$ in \eqref{eq:lambdadef}, so that $\lambda = 1$ when the realization is balanced.
\end{lemma}

\begin{remark} Let us fill in more details of the sketched proof from \cite[Claim A.7]{ECathedral}. We can write each $\dd_s$ (and compositions thereof) as a $Q$-linear combination of elements of $W$, where $Q$ is the fraction field of $R$. For example, when $m_{st} = 3$, we have
\begin{equation} \dd_s \dd_t \dd_s(f) = \frac{(-1)^3 sts(f)}{\alpha_s s(\alpha_t) st(\alpha_s)} + \ldots, \quad \dd_t \dd_s \dd_t(f) =\frac{(-1)^3 sts(f)}{\alpha_t t(\alpha_s) ts(\alpha_t)} + \ldots.\end{equation}
More generally, the coefficient of $w_{s,t}$ in $\dd_s \dd_t \cdots$ (the leading term) will be $(-1)^{m_{st}}$ times the reciprocal of the product of the $s$-aligned choice of roots as in \cite[Definition A.6]{ECathedral}, whereas for $\dd_t \dd_s \cdots$ one instead uses the $t$-aligned choice of roots. One can compare these leading coefficients using \eqref{eq:lambdadef}, and verify that \eqref{eq:rescaled braid} holds for the leading terms. With more combinatorial work one could verify this for all terms. \end{remark}

\begin{remark}
When $W_{s,t}$ acts faithfully and the Chevalley-Shephard-Todd theorem holds, one can bypass this combinatorial work by arguing that any two $R^{s,t}$-linear operators of degree $-2m$ are colinear, and that elements of $W_{s,t}$ are linearly independent over $Q$.

Also, note that both sides of \eqref{eq:rescaled braid} are zero when the action of $W_{s,t}$ is not faithful. If $W_{s,t}$ acts trivially, all Demazure operators are zero. Otherwise the action will factor through a smaller dihedral group, where \eqref{eq:rescaled braid} holds by induction, and then $\dd_s \dd_s=0$ implies the statement.
\end{remark}


% \begin{example} For the usual permutation realization of $S_3$, with simple reflections $\{s,t\}$, we have $\dd_s \dd_t \dd_s = \dd_t \dd_s \dd_t$. There is a different realization of $S_3$, with the same underlying representation, but where $\alpha_t$ and $\alpha_t^\vee$ are multiplied by $\zeta$ and $\zeta^{-1}$ respectively. Then $\dd_s \dd_t \dd_s = \zeta \dd_t \dd_s \dd_t$. \end{example}

% We may use these Demazure operators associated to simple reflections to define the Demazure operator associated to any expression of an element $w \in W$.
% We exclusively use the notation $\un{w}$ to refer to an expression $(s_1,\ldots, s_d)$ of an element $w=s_1\ldots s_d \in W$.

\begin{definition}\label{defn-demazure operator associated to an expression}
For an expression $\un{w}=(s_1,\ldots, s_d)$ of $w\in W$, we define the Demazure operator $\dd_{\un{w}}$ associated to it by setting
$$\dd_{\un{w}}:= \dd_{s_1}\ldots \dd_{s_d}.$$
\end{definition}
When the expression $\un{w}$ is not reduced, \cref{lemma demazure operators coxeter relations} implies that $\dd_{\un{w}}$ is 0. \cref{lemma demazure operators coxeter relations}(2),(3)  together with Matsumoto's theorem implies that  $\dd_{\un{w}}$ is independent of the particular choice of a reduced expression $\un{w}$ of $w$ up to a unit in $\kk.$ We pause to note a conjecture, which to our knowledge is not in the literature at this level of generality (when the braid relations hold only up to scalar).

\begin{conjecture} If $W$ acts faithfully and $\un{w}$ is a reduced expression, then $\dd_{\un{w}}$ is nonzero. \end{conjecture}

We now record how Demazure operators behave for the deformed affine realization. Until the end of this subsection, $(W,S)$ is the affine Weyl group, where $S$ is in bijection with $\Omega = \Z/n\Z$, and $R_{\zeta}$ denotes the graded ring $\Sym_{\Z[\zeta,\zeta^{-1}]}(\Lambda_{\zeta})$ with $\Lambda_{\zeta}$ in degree 2.

% \BE{better first paragraph, referring to earlier conventions}

% Throughout this subsection, $(W, S)=(\tilde{S}_n, \set{s_0, s_1, \ldots, s_{n-1}})$ and we freely identify $S$ with $\Om=\Z/n\Z$ in the obvious way. 
% We often abuse notation and refer to an element of $\Om$ using integer representatives. For example, $n+1$ and $1$ both refer to the same element in $\Om$. 
% We start with some preliminary computations regarding the Demazure operators for the deformed affine realization $\Lambda_{\zeta}$ over $\Z[\zeta,\zeta^{-1}]$ (see \cref{defn-deformed affine realization}). 
% Throughout this section, $R$ denotes the graded ring $\Sym_{\Z[\zeta,\zeta^{-1}]}(\Lambda_{\zeta})$ with $\Lambda_{\zeta}$ in degree 2.

\begin{lemma}\label{lemma demazure operators computation deformed affine}
For $i \in \Om$, let $\dd_i$ denote the corresponding Demazure operator on $R_{\zeta}$.
\begin{enumerate}
    \item For $i \in \Om$, $\dd_i(x_j)=\begin{cases}
        1 &\text{ if $j=i$,}\\
        -\zeta^{-1} &\text{ if $j=i+1$,}\\
        0 &\text{ otherwise.}
    \end{cases}$
    \item For $i,j \in \Om$ such that $s_is_j=s_js_i$, $\dd_i \dd_j= \dd_j \dd_i$.
    \item For $i \in \Om$, $\zeta \dd_{i}\dd_{i+1}\dd_{i}=\dd_{i+1}\dd_i\dd_{i+1}.$
    % \item For $w \in W$, the Demazure operator $\dd_{\un{w}}$ associated to a reduced expression $\un{w}$ of $w$ is independent of the particular choice of the reduced expression up to a unit in $\Z[\zeta,\zeta^{-1}]$
    \end{enumerate}
\end{lemma}
\begin{proof}
    This is straightforward from \cref{lemma demazure operators coxeter relations}. The most interesting part is (3), where the unit $\zeta = \lambda$ agrees with \eqref{eq:lambdadef} when $s = s_i$ and $t = s_{i+1}$, given that
    \begin{equation} \zeta s_i(\alpha_{i+1}) = s_{i+1}(\alpha_i). \end{equation}
    Alternatively, one can compute the unit directly by evaluating both sides on $x_i^2x_{i+1}$.

    % Part (1) is an easy check.
    % Part (2) may be seen as a consequence of \cref{lemma demazure operators coxeter relations}(2), but can also be easily checked directly.
    % Part (3) follows from \cref{lemma demazure operators coxeter relations}(3), where the unit is computed to be $\zeta^{-1}$ by evaluating both sides on $x_i^2x_{i+1}$.
    % One may also do this directly, as outlined in the proof of \cite[Lemma 4.6]{EJY1}, or by using the same arguments as in \cref{lemma- permutation realization from deformed affine} and \cref{lemma-demazure deformed affine vs permutation}.
    % Part (4) is a consequence of (2), (3) and Matsumoto's theorem.
\end{proof}

% \PT{Will comment out the lemma below for the draft.}
% \begin{lemma}\label{lemma deformed nilCoxeter}
%         The algebra $NC(\zeta,n)$ generated by the $\dd_{s_i}$ in End$_{\Z[\zeta,\zeta^{-1}]}(R)$ has the explicit presentation (as an algebra over $\Z[\zeta,\zeta^{-1}]$):
%     Generators $\{D_{s_i} | i \in \mathbb{Z}/n\mathbb{Z}\}$ with relations $D_{s_i}^2=0, \;\;D_{s_i}D_{s_{i+1}}D_{s_i}= \zeta^{-1}D_{s_{i+1}}D_{s_i}D_{s_{i+1}}, \;\; D_{s_i}D_{s_j}=D_{s_j}D_{s_i}$ whenever $s_is_j=s_js_i$.
% \end{lemma}

% \begin{proof}
%     That these relations hold is an immediate check. Let $\mathcal{N}$ denote the algebra defined by the given generators and relations.
%     $\{D_{\underline{w}}\}$, as we run over reduced expressions $\un{w}$ for a given $w\in W$, and over all $w \in W$, forms a spanning set (over $\Z[\zeta,\zeta^{-1}]$) for $\mc{N}$ (and are 0 when the expression is not reduced). 
%     Using the relations and Matsumoto's theorem, we have that for a given $w \in W$, the $D_{\un{w}}$ are independent of the choice of the reduced expression $\un{w}$ of $w$ up to scaling by a unit $\zeta^l$ (for some $l\in \Z$) in $\Z[\zeta,\zeta^{-1}]^\times.$ 
%     Let's fix a choice of reduced expression $\un{w}$ for each $w \in W$, and define $D_w$ to be $D_{\un{w}}$, so that $\{D_w\}_{w \in W}$ forms a spanning set for the algebra $\mc{N}$. 
%     We now proceed to show that $\{D_{w}\}_{w \in W}$ is a linearly independent set and that the specialization map $\phi:\mc{N}\rightarrow \End_{\Z[\zeta,\zeta^{-1}]}(R): D_{s_i} \rightarrow \dd_{s_i}$ is an isomorphism. 
%     Indeed, both these statements would follow if we show that the $\{\dd_{w}:= \dd_{\un{w}}= \phi(D_w)\}_{w \in W}$ is a linearly independent set (and hence a basis of $NC(\zeta,n)$).
%     \begin{claim}
%         $\{\dd_w\}_{w \in W}$ is a $\Z[\zeta,\zeta^{-1}]$-linearly independent set of elements in $NC(\zeta,n)$.
%     \end{claim}
%     \begin{proof}[Proof of Claim]
%         Since we have a grading on $R$ with respect to which the $\dd_{s_i}$ are operators of degree $-2$, we have that the $\dd_w$ for $w\in W$ are of degree $\ell(w)$. 
%         Hence it suffices to show linear independence of the $\dd_w$ for $w \in W$ of a fixed length $l.$
%         \PT{Aaaah.... This is harder than I thought, but is also not necessary for this paper... Maybe I'll just stick with saying that the deformed braid relations hold....}
%     \end{proof}    
%     The lemma then follows as explained in the paragraph preceding the claim.
%     \end{proof}


% \subsubsection{Symmetry relative to the Rotation Operator $\sigma$}\label{subsubsec rotational symmetry for deformed affine realization}

Recall that we had defined the rotation operator $\sigma: \Omega \rightarrow \Omega: i \mapsto i+1$  and extended it to an operator on $\Lambda_{\zeta}$ such that $\sigma(\zeta)=\zeta$ and $\sigma(x_i)=x_{\sigma(i)}=x_{i+1}$ (see \cref{defn-rotation operator}). Hence $\sigma$ naturally extends to a ring automorphism $\sigma: R\xrightarrow{\sim} R$ which preserves the grading.

\begin{lemma} We have
\begin{equation} \label{eq rotational symmetry deformed affine} \sigma \circ \dd_i \circ \sigma^{-1} (f) = \dd_{\sigma(i)}(f). \end{equation}
\end{lemma}

\begin{proof} One has $\sigma(s_{k}(x_{j}))=s_{k+1}(x_{j+1})=s_{\sigma(k)}(\sigma(x_j))$ for any $k, j\in \Om$, so that $\sigma\circ s_k \circ \sigma^{-1}=s_{\sigma(k)}$ as operators on $R$.
It follows that
\[ \sigma\dd_i \sigma^{-1}(f)=\sigma \left(\frac{\sigma^{-1}(f)-s_i(\sigma^{-1}(f))}{x_i-\zeta x_{i+1}} \right)
=\frac{f-s_{i+1}(f)}{x_{i+1}-\zeta x_{i+2}}=\dd_{\sigma(i)} (f). \]
\end{proof}
% so that $\sigma \circ \dd_i \circ \sigma^{-1} = \dd_{\sigma(i)}$. 
% A similar relation also holds for all powers of $\sigma$.
We refer to this as \emph{rotational symmetry} or \emph{color symmetry} in the later sections.

The action of $\sigma$ on $S=\Om$ also extends naturally to an action on the set of words in $S$ and on $W$. 
Given an expression $\un{w}:=(s_{i_1}, \ldots s_{i_k})$ for $w=s_{i_1}\ldots ,s_{i_d}$, we have $\sigma(\un{w}):=(s_{\sigma(i_1)}, \ldots ,s_{\sigma(i_d)})$ as an expression for $\sigma(w):=s_{\sigma(i_1)}\ldots s_{\sigma(i_d)}$. 
Then \cref{eq rotational symmetry deformed affine} implies that 
\begin{equation}
    \sigma\dd_{\un{w}}=\dd_{\sigma(\un{w})}\sigma.
\end{equation}


\begin{remark}\label{rmk symmetry for rotation operator}
    In the special case when $f \in R$ is such that $\dd_{\un{w}}(f)$ is of degree 0, we have that $\sigma(\dd_{\un{w}}(f))=\dd_{\un{w}}(f)=\dd_{\sigma(\un{w})}(\sigma(f)).$
\end{remark}

%Commented out the original remark below and added it as a subsubsection
\comm{
\begin{remark}(Symmetry for the Frobenius realization with respect to the rotation operator $\sigma$)\label{rmk symmetry for rotation operator}
In the proof, we used a certain symmetry within our Frobenius realization with respect to the rotation operator $\sigma$ (defined in \cref{defn-rotation operator}), which we now elaborate.
Recall that we had defined the rotation operator to be $\sigma: \Om \rightarrow \Om: i \mapsto i+1$.
Hence $\sigma$ naturally extends to an automorphism $\sigma: R\xrightarrow{\sim} R$ such that $\sigma(x_i)=x_{\sigma(i)}=x_{i+1}$ for all $i \in \Om$.
Then it follows that $\sigma(s_{k}(x_{j}))=s_{k+1}(x_{j+1})=s_{\sigma(k)}(\sigma(x_j))$ for any $k, j\in \Om$, so that $\sigma\circ s_k \circ \sigma^{-1}=s_{\sigma(k)}$ as operators on $R$.
It further follows that
\[\sigma\dd_i \sigma^{-1}(f)=\sigma \left(\frac{\sigma^{-1}(f)-s_i(\sigma^{-1}(f))}{x_i-\zeta x_{i+1}} \right)
=\frac{f-s_{i+1}(f)}{x_{i+1}-\zeta x_{i+2}}=\dd_{\sigma(i)} (f)\]
so that $\sigma \circ \dd_i \circ \sigma^{-1} = \dd_{\sigma(i)}$. 
A similar relation also holds for all powers of $\sigma$.
We may sometimes refer to this as \emph{rotational symmetry} or \emph{color symmetry} in the later sections.
As an example, in the proof of \cref{lemma dual basis left ext}, we used this symmetry to simplify \eqref{eq 4 dual basis left ext}:
\begin{align*}\dd_k\ldots \dd_2\left(\zeta^{k-1}e_{k-1}(\zeta x_3, \ldots, \zeta^{k-1} x_{k+1})\right)&=\sigma \left(\dd_{k-1}\ldots \dd_1\left(\zeta^{k-1}e_{k-1}(\zeta x_2, \ldots, \zeta^{k-1} x_{k})\right)\right).
\end{align*}
Since $\sigma$ acts trivially on the degree 0 elements of $R$, we get \eqref{eq 5 dual basis left ext}.
\end{remark}
}



% \subsubsection{Symmetry relative to the Reflection Operator $\tau$}\label{subsubsec reflection symmetry for reflection operator}

Recall that we had defined the reflection operator $\tau: \Om \rightarrow \Om: k \mapsto -k$, and extended it to an operator on $\Lambda_{\zeta}$ 
such that $\tau(\zeta)=\zeta^{-1}$ and $\tau(x_i)=x_{n-i+1}$ (see \cref{defn-reflection operator}). This extends to a ring automorphism $\tau:R \xrightarrow{\sim}R$.

\begin{lemma} We have
\begin{equation} \label{eq:reflection symmetry for reflection operator} \tau \circ \dd_i \circ \tau^{-1}(f) = - \zeta \dd_{\tau(i)}(f). \end{equation}
\end{lemma}

\begin{proof}
It is straightforward to check that $\tau(s_k(x_i))=s_{n-k}(\tau(x_i))$ for any $i, k\in \Om$, i.e., $\tau\circ s_k\circ\tau^{-1}=s_{\tau(k)}$.
It follows then that 
\[\tau\dd_i \tau^{-1}(f)= \tau \left( \frac{\tau^{-1}(f)-s_{i}(\tau^{-1}(f))}{x_i-\zeta x_{i+1}}
\right) = \frac{f-s_{n-i}(f)}{x_{n-i+1}-\zeta^{-1}x_{n-i}}=-\zeta \dd_{n-i}(f). \] 
\end{proof}
% so that $\tau\circ\dd_i \circ \tau^{-1}=-\zeta \dd_{\tau(i)}$ or $\tau \circ \dd_i=-\zeta \dd_{\tau(i)}\circ \tau$.

The action of $\tau$ on $S=\Om$ extends naturally to an action on the set of words in $S$, and on $W$. 
Given an expression $\un{w}:=(s_{i_1}, \ldots s_{i_d})$ for $w=s_{i_1}\ldots s_{i_d}$, we have $\tau(\un{w}):=(s_{\tau(i_1)}, \ldots s_{\tau(i_d)})$ as an expression for $\tau(w):=s_{\tau(i_1)}\ldots s_{\tau(i_d)}$. 
Then we have that 
\begin{equation}
    \tau\dd_{\un{w}}=(-\zeta)^{d}\dd_{\tau(\un{w})}\tau.
\end{equation}


%==========================================================================
\subsection{Frobenius hypercube realizations and the deformed affine setting}
\label{subsec-frob realizations}
\label{subsec-frob realization deformed affine}
%==========================================================================

The representation of $W$ produces the ring $R$; the additional data provided by a realization (roots and coroots) pins down the Frobenius trace $R \to R^s$. Let us pin down all the Frobenius traces.

\begin{defn}\label{defn frobenius hypercube realization} A \emph{Frobenius (hypercube) realization} or \emph{Frealization} is a realization equipped with a choice of a compatible partial Frobenius hypercube, with $\Gamma = S$, $\mathcal{F}$ the collection of finitary subsets of $S$, and $R^I$ the subring of $W_I$-invariants in $R$ for $I\in \mc{F}$.
In other words, we have a
Frobenius structure $\dd^J_I$ on $R^I \subset R^J$ for each finitary pair $J \subset I \subset
S$ such that $\dd^J_I \dd^K_J = \dd^K_I$ whenever $K \subset J$.
(In particular, we assume that such a Frobenius structure exists.) 
We also require that 
\begin{itemize}
\item The Frobenius trace $R \to R^s$ must simply be $\dd_s$, as determined by the realization. 
\item The Frobenius trace $\dd_I : R \to R^I$ must agree with $\dd_{\un{w}_I}$, up to a unit in $\kk$, for some reduced expression $\un{w}_I$ of $w_I$.
\end{itemize}
\end{defn}

\begin{remark} The second condition, that $\dd_I$ and $\dd_{\un{w}_I}$ agree up to unit, is automatic by \cref{lem Frobenius structures unique up to unit} so long as $\dd_{\un{w}_I}$ is nondegenerate. A prerequisite for $\dd_{\un{w}_I}$ to be nondegenerate is for it to be surjective as a map $R \to R^I$, which is called \emph{generalized Demazure surjectivity} \cite[Definition 3.5]{EKLP}.\end{remark}

\begin{remark} It seems as though the choice of $\dd_s$ is fixed whereas the choice of $\dd_I$ is an invertible scalar by \cref{lem Frobenius structures unique up to unit}. In truth, $\dd_s$ can also be rescaled by rescaling the choice of roots and coroots (i.e. changing the underyling realization). An alternate approach would be to omit the realization from the data, and recover it from the Frobenius extension structure for $R^s \subset R$. \end{remark}

\begin{example}\label{example nice realizations are frobenius}
If the realization $\Lambda$ is balanced (see \S\ref{subsec-balanced}) then we unambiguously define $\dd_w$ to be $\dd_{\un{w}}$ for any reduced expression $\un{w}$ of $w$. 
In this case, we can then define $\dd^\emptyset_I:= \dd_{w_I}$ for a finitary $I\subset S$, and 
$$\dd^J_I:=\dd_{w_Iw_J^{-1}}$$
for $J \subset I \subset S$. 
Compatibility will hold. If the realization is also faithful and satisfies generalized Demazure surjectivity, then $\dd^J_I$ is nondegenerate, and this data yields a Frealization. This is proven in \cite[Theorem 4.3]{EKLP} under these exact assumptions on the realization (see \cite[Section 3.1]{EKLP}).
\end{example}


% The deformed affine realization that we want to work with is not a balanced realization- for example, we have $\dd_{s_1}\dd_{s_2}\dd_{s_1}=\zeta^{-1}\dd_{s_2}\dd_{s_1}\dd_{s_1}$ (see \cref{lemma demazure operators computation deformed affine}). 

% But the $\dd_{s_i}$, $i\in \Om=\Z/n\Z$ still satisfy the braid relations up to a unit in the ring $\kk$, so we have that $\dd_{\un{w}}\in\kk^\times \dd_{\un{w}'}$ if $w=w' \in W$. 
% So we may still define $\dd_I^\emptyset$ to be an element of $\kk^\times\dd_{\un{w}_I}$ satisfying some normalization condition, for any reduced expression $\un{w}_I$ for $w_I.$ 
% One can also define $\dd^J_I$  which satisfies $\dd_I^J\dd^K_J=\dd^K_J$ to obtain a Frobenius realization. 
% See \cref{defn deformed affine frobenius realization} for the details.

% \comm{
% \begin{remark}
%     A similar procedure can be done to obtain a Frobenius realization $\Lambda$ from any realization over a commutative domain $A$ where the $\dd_{s}, s\in S$ satisfy the braid relations up to a unit in $A$ with some additional assumptions. \PT{I think I know the exact things needed, but skipping for now. Maybe I'll just cite Abe.}    
% \end{remark}
% }

Our goal for the rest of this section is to define a Frobenius realization extending the deformed affine realization. Once again, $(W,S)$ is the type $A$ affine Weyl group with our standard conventions.

\begin{defn}
We say that a nonempty subset $I \subset S$ is \emph{connected} if $I=\set{i, i+1, \ldots i+k}$ for some integers $i$ and $k$ (considered modulo $n$). If $I\subset S$ is not connected, we say it is \emph{disconnected}. 
A \emph{connected component} of a disconnected subset is a maximal connected subset.
\end{defn}

\begin{defn}\label{defn deformed affine frobenius realization}
We make the following choices for the deformed affine realization $\Lambda_{\zeta}$ (see Definition \ref{defn-deformed affine realization}) defined over $\Z[\zeta,\zeta^{-1}]$, letting $R = R_{\zeta}$.
\begin{itemize}
    \item We fix a system of \textbf{positive roots} $\Phi_I^+$ for each subset $I\subset S$. 
    \\If $I=\{s_1, s_2, \ldots , s_{d-1}\}$, we choose our positive roots to be 
    $$x_i - \zeta^{j-i}x_j \hspace{20pt} \text{for } 1\leq i<j\leq d.$$
    If $I$ is connected and $I = \sigma^l(J)$ for some $l \in \Z$, then we set $\Phi^+_I = \sigma^l(\Phi^+_J)$.
    % If $I$ is some other connected subset of $S$, it is related to the previous set by some $\sigma^l$ for some $l \in \Z$ and we define the positive roots to be $\sigma^l$ acting on the positive roots in the previous case.
    \\ For $I$ disconnected, we take a disjoint union of the positive roots for its connected components.
    \\ Note that for $I \subset J$ finitary, $\Phi^+_I \subset \Phi^+_J$.
    \item Now for finitary $I \subset S$, we define $\mu_I$ to be 
    $$\mu_I:= \prod_{\al \in \Phi_I^+} \al.$$
    We will use these $\mu_I$ when defining the Frobenius trace map $\dd_I$ so that the $\mu_I$ coincide with the product-coproduct element of the Frobenius extension $R^I \hookrightarrow R$ thus defined.
    We also define $$\mu_I^J:= \prod_{\al \in \Phi_I^+\setminus \Phi_J^+}\al\;\;,$$
     and this will coincide with the product-coproduct element for the Frobenius extension $R^I \hookrightarrow R^J$. 
     Note that for $K\subset  J \subset I$, we have that $\mu^K_I=\mu^J_I\mu^K_J.$
    \item Now we fix Frobenius trace maps as follows.
    \\For $I \subset S$ finitary and any reduced expression $\un{w}_I$ of the longest element $w_I$, we define $\dd_I$ to be the (unique) rescaling of $\dd_{\un{w}_I}$ so that $\dd_I(\mu_I)= |W_I|.$  Moreover, for $J \subset I \subset S$, we define $\dd_I^J$ to be the unique rescaling of $\dd_{\un{w_Iw_J^{-1}}}$ (for any reduced expression $\un{w_Iw_J^{-1}}$ of $w_Iw_J^{-1}$) such that $\dd^J_I(\mu_I^J)=|W_I|/|W_J|$.
    
    % $\dd_{\un{w}_I}$ gives a Frobenius trace map $R\rightarrow R^I$. 
    % These are independent of the choice of reduced expression up to a unit in $\Z[\zeta,\zeta^{-1}]$ (see lemma \ref{lemma deformed nilCoxeter} below), and we define $\dd_I$ to be the (unique) rescaling of the trace map so that $\dd_I(\mu_I)= |W_I|,$ where $\mu_I$ is as defined above.
    % Moreover, for $J \subset I \subset S$, we define $\dd_I^J$ to be a rescaling of $\dd_{\un{w_Iw_J^{-1}}}$ for any reduced expression $\un{w_Iw_J^{-1}}$ of $w_Iw_J^{-1}$ such that $\dd^J_I(\mu_I^J)=|W_I|/|W_J|$.
    % \PT{@BE I know that I'm using slightly different notation $\un{w}_I$ vs $\un{w_I}$ for the case $I=\emptyset$. I like the former better, but I couldn't think of a way to do a similar thing for $w_Iw_J^{-1}$. Maybe define $w^J_I:=w_Iw_J$ and then write $\un{w}^J_I$? }
\end{itemize}
\end{defn}

In fact, in \cref{subsec:stdrex} we will choose a particular reduced expression $\un{w}_I^{\std}$ such that $\dd_I = \dd_{\un{w}_I^{\std}}$ on the nose.

\begin{lemma}\label{lemma deformed affine frobenius realization is frobenius} The above definition yields a well-defined Frealization. \end{lemma}

\begin{proof} 
Fix $k \in \Omega$. Combining \cref{lemma-deformed affine restricted to Sn} with the symmetry $\sigma$, we can identify the restriction of the deformed affine realization to $W_{\hh{k}}$ with the permutation representation  $V_{\perm}=\bigoplus_{i=1}^n \Z[\zeta,\zeta^{-1}] y_i$ of $S_n$ via 
\begin{equation}\label{eq identifying deformed affine restricted to W-k-hat with permutation}
y_1=x_{k+1}, y_2=\zeta x_{k+2}, \ldots ,y_{n-1}=\zeta^{n-2} x_{k-1}, y_n=\zeta^{n-1} x_k.     
\end{equation}
Through this identification, the simple reflections $s_{k+i}$ (for $1 \leq i \leq n-1)$ act as simple transpositions $t_i$ swapping $y_i$ and $y_{i+1}$.
The permutation realization has simple roots $\alpha_{t_i} = y_i-y_{i+1}$. The identification does not preserve the roots and coroots, and we have
\begin{equation} \alpha_{k+i} \mapsto \zeta^{-(i-1)} \alpha_{t_i}, \quad \dd_{k+i} \mapsto \zeta^{i-1} \dd_{t_i}. \end{equation}
The other roots in $\Phi^+_{\hh{k}}$ are also sent to scalar multiples of positive roots in $\{y_i - y_j\}_{i > j}$. To disambiguate from $\dd_{k+i}$, we write $\tilde{\dd}_{k+i}$ for $\dd_{t_i}$.

Let $I \subset \hat{k}$ be any parabolic subgroup, and let $\un{w}_I$ be any reduced expression for $w_I$. The permutation realization (over $\Z$ or $\Z[\zeta, \zeta^{-1}]$) is a balanced realization and satisfies generalized Demazure surjectivity, so the discussion of \cref{example nice realizations are frobenius} applies. Thus the composition $\tilde{\dd}_{\un{w}_I}$ is a non-degenerate map $R \to R^I$, and it is well-known (see \cite{Demazure}) that the corresponding product-coproduct element is the product $\tilde{\mu}_I$ of the positive roots for $I$ (a subset of $\{y_i - y_j\}_{i < j}$). Since $\dd_{\un{w}_I}$ agrees with $\tilde{\dd}_{\un{w}_I}$ up to a power of $\zeta$, then $\dd_{\un{w}_I}$ is also nondegenerate. By \cref{lem Frobenius structures unique up to unit}, its product-coproduct element agrees up to a power of $\zeta$ with $\tilde{\mu}_I$, and thus also with $\mu_I$.

Consequently, $\dd_{\un{w}_I}(\mu_I) = \zeta^l |W|$ for some $l \in \Z$. We can define $\dd_I$ as above, and it will be non-degenerate. Its product-coproduct element will be a scalar multiple of $\mu_I$, and the normalization condition $\dd_I(\mu_I) = |W_I|$ implies that it is $\mu_I$ on the nose. Moreover, $\dd_I$ is independent of the choice of $\un{w}_I$, again by \cref{lem Frobenius structures unique up to unit}. Similar statements can be made about $J \subset I \subset \hat{k}$ and the relative longest element $w_Iw_J^{-1}$.

We check the other requirements of \cref{defn frobenius hypercube realization}. Since $\mu_s = \al_s$, we have  $\dd^{\emptyset}_s =\dd_s$ as required. The compatibility condition \eqref{eq:mumult2} follows from the definition. \end{proof}

% , and the corresponding Demazure operators given by

% \begin{equation}
% \dd_{t_i}(f)=\displaystyle \frac{f-t_i(f)}{y_i-y_{i+1}}=\frac{f-s_{k+i}(f)}{\zeta^{i-1}(x_{k+i}-\zeta x_{k_i+1})}=\zeta^{-(i-1)}\dd_{k+i}=:\tilde{\dd}_{k+i}.\label{eq demazure deformed affine vs permutation}
% \end{equation}

% By \cref{example nice realizations are frobenius}, the Demazure operators $\tilde{\dd}_{k+i}$ satisfy the braid relations and we may define $\tilde{\dd}_{w_I}=\dd_{\un{w}_I}$ for any reduced expression $\un{w}_I$ of $w_I$, for $I \subset \hh{k}$.
% By \cref{eq demazure deformed affine vs permutation}, $\tilde{\dd}_{\un{w}_I}$ coincides with $\dd_{\un{w}_I}$ up to a unit of the form $\zeta^l$. 
% It is well known that the permutation realization is faithful and satisfies generalized Demazure surjectivity, so by \cref{example nice realizations are frobenius}, we have that $\tilde{\dd}_I$, and hence $\dd_{\un{w}_I}$ are non-degenerate for $J\subset I \subset \hh{k}$.
% For $J \subset I \subset \hh{k}$, one uses a similar argument to conclude that $\dd_{\un{w_Iw_J^{-1}}}$ is non-degenerate as well.

% Hence we have that the $\dd_I, \dd_I^J$ are Frobenius trace maps for $J\subset I \subset \hh{k}$, and that the product-coproduct elements coincide with those associated to $\tilde{\dd}_I, \tilde{\dd}_I^J$ up to a unit.
% The product-coproduct element $\tilde{\mu}_I$ for the permutation realization is known to be a product of positive roots (given by $y_i-y_j, i<j$ for the transposition $(i,j)$) corresponding to $I$.
% One can check that the identification in \cref{eq identifying deformed affine restricted to W-k-hat with permutation} takes  $\mu_I$ (in \cref{defn q-deformed affine Frobenius realization}) to $\tilde{\mu}_I$, up to a unit of the form $\zeta^l$ for $l \in \Z$.
% Hence $\mu_I$ coincides with the product-coproduct element up to a unit, and by the normalization condition for $\dd_I$ in \cref{defn q-deformed affine Frobenius realization}, they have to be equal.
% For $J \subset I \subset \hh{k}$, a similar argument along with the fact that $\mu_I^J=\mu_I/\mu_J$ shows that $\mu_I^J$ in \cref{defn deformed affine frobenius realization} equals the product-coproduct element for the Frobenius extension $R^I \hookrightarrow R^J$ with trace $\dd_I^J$.

% Since every finitary subset of $S$ is in some maximal finitary subset, for arbitrary finitary subsets $J \subset I \subset S$, we have that $\dd_I^J$ is a Frobenius trace with product-coproduct element $\mu_I^J$.

% Now we check the additional compatibility conditions in \cref{defn frobenius hypercube realization}. \PT{End of addition.}\BE{added this} Since $\mu_s = \al_s$, we have  $\dd^{\emptyset}_s =\dd_s$ as required. For $K \subset J \subset I$, one can check that $\dd^J_I\dd^K_J=\dd^K_I$ (using the fact that $\mu^K_I=\mu^J_I\mu^K_J$). \PT{Commented out the last two lines.}
%It remains to show that one can choose a reduced expression $\un{w}_I$ for which $\dd_{\un{w}_I}$ is a valid Frobenius trace, and whose product-coproduct element is an invertible multiple of $\mu_I$, see \BE{ref remark}, and similarly for each $I \subset J$. 
%We postpone this to the next section where we find explicit dual bases, see \cref{corollary- deformed affine is a frobenius realization}.

% Each finitary $I$ is contained in a maximal parabolic subgroup $\hh{k}$ for some $k \in \Omega$. The restriction of $\Lambda_{\zeta}$ to $W_{\hh{k}}$ is isomorphic to the permutation representation of $W_{\hh{k}} \cong S_n$ (see \BE{ref}). Since this isomorphism sends the positive roots $\Phi_I$ to scalar multiples of the ordinary positive roots, it also sends $\dd_i$ to a scalar multiple of the ordinary Demazure operator for $S_n$. Consequently, one can deduce the desired statement from the fact that $\Z[y_1, \ldots, y_n]^{S_n} \subset \Z[y_1, \ldots, y_n]$ is a Frobenius extension with its usual trace and coproduct-product element.




\begin{remark}\label{rmk dual basis from permutation realization} 
The identification \cref{eq identifying deformed affine restricted to W-k-hat with permutation} may be used for constructing dual basis for the Frobenius extensions in the deformed affine realization, by rescaling from the permutation realization. In \cref{subsec-dual basis} we pick dual bases explicitly.
\end{remark}

{
% \begin{lemma}
% The Demazure operators from the permutation realization $\dd_{t_i}$ correspond to
% $\zeta^{-(i-1)}\dd_{s_{k+i}}.$
% \end{lemma}

% \begin{corollary}\label{corollary Frobenius trace in deformed affine vs permutation}
%     For $J \subset I \subset \hh{k}$, the trace maps $\dd_I, \dd_J$ and $\dd_I^J$ in \cref{defn deformed affine frobenius realization} respectively coincide, up to a unit of the form $\zeta^l$ for some integer $l$, with the corresponding Frobenius trace maps of the permutation realization obtained by restricting to $W_{\hh{k}}$.
% \end{corollary}


% \comm{
% Now, suppose $I\subset S$ is finitary, and is contained in a maximal finitary $\tilde{I}=S \setminus \set{k}$.
% For any $K \subset \tilde{I}$, let $\tilde{\dd}_K$ denote the Frobenius trace map corresponding to the permutation realization obtained by restriction to $(W_{\tilde{I}}, \tilde{I})$. 
% We have seen in the discussion above \cref{corollary- deformed affine is a frobenius realization} that $\tilde{\dd}_I=\zeta^l \dd_I$ for some integer $l$.
% Suppose $\set{e_i}_{i \in W_{\tilde{I}}}, \set{f_j}_{j \in W_{\tilde{I}}}$ are dual bases for the extension $R^I \hookrightarrow\ R$ with respect to the trace $\tilde{\dd}_I$ (there are well-known dual basis for the Frobenius extensions from the permutation realization, for example, see \PT{Figure out the best place to cite this!}).
% Then $\set{\zeta^le_i}$ and $\set{f_j}$ form a dual basis for $\dd_I$. 
% Let $\tilde{\mu}_I$ is the product-coproduct element for the Frobenius extension $R^I \hookrightarrow R$ with trace $\tilde{\dd}_I$ (i.e., this is the product-coproduct element coming from identifying the restriction to $\tilde{I}$ with the permutation realization). 
% Then we have that $\zeta^l\tilde{\mu}_I=\zeta^l \sum e_if_i$ be the product-coproduct element for the same extension $R^I \hookrightarrow R$ but with trace $\dd_I$.
% We may use this observation to prove the following lemma.}


% \begin{proof}
% For a maximal finitary subset $I$, one can check that the identification in \cref{lemma- permutation realization from deformed affine} takes  $\mu_I$ (in \cref{defn q-deformed affine Frobenius realization}) to the product of the positive roots for the parabolic subgroup $W_I$ in the permutation realization, up to a scalar of the form $\zeta^l, l \in \Z$.
% It is well known that for the permutation realization, the product of positive roots of $W_I$ coincides with the product-coproduct element for $R^I \hookrightarrow R$.
% Since we already know that the corresponding Frobenius trace maps are the same up to a scalar, the normalization condition for $\dd_i$ in \cref{defn q-deformed affine Frobenius realization} then shows that $\mu_I$ coincides with the product-coproduct element.

% The statement for arbitrary finitary $I \subset S$ follows by the same argument after choosing a maximal finitary subset containing $I$.
% For  $J\subset I$, the statement regarding $\mu_I^J$ follows from the fact that $\mu_I^J= \dfrac{\mu_I}{\mu_J}$.
% %commeting out old proof
%      \comm{For simplicity, let's assume $I\subset S\setminus\set{0}=\tilde{I}$. 
%      The discussion before the lemma implies that the product-coproduct element is $\zeta^l \tilde{\mu}_I$ where $l$ is such that $\tilde{\dd}_I=\zeta^l \dd_I$.
%      Since $\tilde{\dd}_I$ is the trace from the permutation realization of $S_n$, we know that 
%      \[\tilde{\dd}_I(f)= \frac{\sum_{w \in W_I}(-1)^{\ell(w)}w(f)}{\prod_{\al \in \tilde{\Phi}_I^+}\alpha} ,\]
%      where $\tilde{\Phi}_I^+$ is the collection of positive roots from the permutation realization.
%      To be precise, for a reflection $w=(i,j)$ for $1\leq i<j \leq n$, we have the corresponding positive root $\tilde{\al}_w= y_i-y_j=\zeta^{i-1}x_i-\zeta^{j-1}x_j$, and $\tilde{\Phi}^+_I=\set{\tilde{\al}_t}$ where $t$ runs over reflections in $W_I$.
%      This explicit description makes it clear that $\tilde{\mu}_I=\zeta^m \mu_I$ for some $m \in \Z$. 
%      Since 
%      $\tilde{\dd}_I=\zeta^l \dd_I$, we have that 
%      $$|W_I|=\dd_I(\text{product-coproduct element for 
%      $R^I \hookrightarrow R$})=\zeta^{m-l}\tilde{\dd}_I(\tilde{\mu_I})=\zeta^{m-l}|W_I|,$$
%      so $m=l.$
%      It follows that the product-coproduct element in question coincides with $\mu_I$.
%      For arbitrary $\tilde{I}$, one may use rotational symmetry (i.e., symmetry with respect to $\sigma$) to argue the same.

%     For $J\subset I$, the statement regarding $\mu_I^J$ follows from the fact that $\mu_I^J= \dfrac{\mu_I}{\mu_J}.$}
% %end of commented out old proof.
%\end{proof}
}

Before continuing to explore this Frealization, we check some technical properties which were introduced in \cite{EWSFrob} to ensure proper behavior of its diagrammatic category.

\begin{defn}\cite[Definition 1.7]{EWSFrob}\label{defn-condition star} We say that a Frobenius hypercube satisfies condition $\star$ if, for every $I \ne J \in \mc{F}$ with $I \cup J \in \mc{F}$, setting $K = I \cap J$ and $L = I \cup J$, we may choose a basis $\{x_\alpha\}$ of $R^K$ over $R^I$ such that $x_\alpha \in R^J$. \end{defn}
Note that in \cite{EWSFrob}, this is stated a bit differently, where it is required to have a pair of dual bases for $R^K$ over $R^I$ where one of the bases lives in $R^J$. 
But if we have any basis of $R^K$ over $R^I$ already in $R^J$, then we can find its dual basis by Remark \ref{rmk-non-degeneracy and dual basis}. The point to emphasize is that typically the dual basis will not live within $R^J$.

\begin{defn}\cite[Definition 1.8]{EWSFrob} \label{defn- no mu zero deivisors} We say that the Frobenius hypercube $\Gamma$ has \emph{no $\mu$-zero
    divisors} if $\mu_I \in R$ is not a zero divisor for all $I$ in the hypercube $\Gamma$.
\end{defn}



\begin{defn}\cite[Definition 1.10]{EWSFrob}\label{defn-condition R3} We say that a Frobenius hypercube satisfies condition R3, if for every $I \in \mc{F}$ and distinct  $i, j, k \in \Gamma \setminus I$ with $I \sqcup\set{i,j,k} \in \mc{F}$, we have that ${\mu_{Iij} \mu_{Iik} 
\mu_{Ijk} \mu_I} \mid {\mu_{Iijk} \mu_{Ii} \mu_{Ij}\mu_{Ik}}$.
\end{defn}

\begin{remark}
Note that in \cite{EWSFrob}, the phrasing for condition R3 is slightly different. Translated to our notation, \cite[Definition 1.10]{EWSFrob} asks that
${\mu^I_{Iij} \mu^I_{Iik} 
\mu^I_{Ijk}} \mid {\mu^I_{Iijk} \mu^I_{Ii} \mu^I_{Ij}\mu^I_{Ik}}$,
which is equivalent to the condition mentioned above, when there are no $\mu$-zero divisors.

As mentioned in \cite[Remark 1.11]{EWSFrob}, when $R$ is a UFD and 
 $\mu^I_{Iij}, \mu^I_{Ijk}, \mu^I_{Iik}$
 are relatively prime, the condition R3 automatically holds. 
\end{remark}

\begin{lemma}\label{lemma deformed affine frobenius satisfies extra conditions}
   The deformed affine realization equipped with the Frobenius structure in \cref{defn deformed affine frobenius realization} satsfies the conditions in 
   \cref{defn-condition star}, \cref{defn- no mu zero deivisors}, \cref{defn-condition R3}.
\end{lemma}

\begin{proof} Clearly this Frealization has no $\mu$-zero divisors, since each ring $R^I$ is a domain and each $\mu_I$ is nonzero. Given that each $\mu_I$ is defined as a product of a certain set of roots, condition R3 follows immediately from the inclusion-exclusion principle.

For condition $\star$, we use similar arguments as in the proof of \cref{lemma deformed affine frobenius realization is frobenius}, passing to the permutation realization by restricting to a maximal finitary parabolic subset containing $I \cup J$. The property $\star$ for the permutation realization over any commutative domain follows from \cite[Theorem 4.15, Example 4.17]{EKLP}. The assumptions for \cite[Theorem 4.15]{EKLP} are that the realization is faithful, balanced, and satisfies generalized Demazure surjectivity (see \cite[Section 3.1]{EKLP}), all of which are known to hold for the permutation realization.

\end{proof}


% It is easy to show that the conditions in  \cref{defn- no mu zero deivisors}, \cref{defn-condition R3} are satisfied for the permutation realization of $S_n$ over any commutative ring, using the fact that the positive roots are of the form $y_i-y_j, i<j$. 
% Showing \cref{defn-condition star} for the permutation realization over any commutative domain is harder, but follows from \cite[Theorem 4.15, Example 4.17]{EKLP24}. The assumptions for \cite[Theorem 4.15]{EKLP24} are that the realization is faithful, balanced, and satisfies generalized Demazure surjectivity (see \cite[Section 3.1]{EKLP24}), all of which are known to hold for the permutation realization.
% The following lemma then follows from similar arguments as in the proof of \cref{lemma deformed affine frobenius realization is frobenius}, after passing to the permutation realization by restricting to maximal finitary parabolic subsets.





\subsection{Choosing reduced expressions} \label{subsec:stdrex}

Now we find the promised reduced expression for $w_I$, for which the corresponding operator $\dd_{\underline{w}_I}$ agrees with $\dd_I$ on the nose. However, when $J \subset I$ it will typically be impossible to find a reduced expression for $w_I w_J^{-1}$ for which the corresponding operator agrees with $\dd^J_I$ on the nose, see \cref{rmk:wontworkforIJ}.


%  For $J \subset I$ both finitary, the fact that the extensions $R^I \hookrightarrow R^J$ is a Frobenius extension with the trace $\dd_I^J$ as defined above will be shown later (\cref{corollary- deformed affine is a frobenius realization}).
%  Hence we can think of the deformed affine realization as a graded Frobenius realization, and obtain a (partial) graded Frobenius hypercube from it.


% Note that in the definition of $\dd_I^J$ in \cref{defn deformed affine frobenius realization} , we scale $\dd_{\un{w_Iw_J^{-1}}}$ by a unit in $\Z[\zeta,\zeta^{-1}]$.
% The following lemma and \cref{corollary-normalization} actually shows that for any $I \subset S$ finitary, the unit can be chosen to be 1 if we choose the ``nicest'' reduced expression for $w_I$: 

\begin{lemma}\label{lemma- demazure standard rex}
    We have $\dd_{\{1,2, \ldots ,k\}}=\dd_1\dd_2\ldots \dd_k \dd_{\{1,2,\ldots , k-1\}}$ for any $k<n.$ \\Hence for the reduced expression
    $$\un{w}^{\std}_{\{1,2, \ldots,k\}}:=(s_1, s_2, \ldots, s_k, s_1,s_2, \ldots , s_{k-1}, \ldots, s_1,s_2,s_1)$$
    of $w_{\{1,2, \ldots,k\}}$, we have that
    $\dd_{\{1,2, \ldots ,k\}}=\dd_{\un{w}^{\std}_{\{1,2, \ldots,k\}}}$.
\end{lemma}
\begin{proof}
    %I have this fully written down in my tablet, induction.
    % Since the Frobenius trace map we defined coincide with $\dd_{\un{w}_{\{1,2, \ldots,k\}}}$  upto a unit in $\Z[\zeta, \zeta^{-1}]$,
    It suffices to show that $\dd_1\dd_2\ldots \dd_k \dd_{\{1,2,\ldots , k-1\}}$ satisfies the normalization condition, i.e., 
    $$\dd_1\dd_2\ldots \dd_k \dd_{\{1,2,\ldots , k-1\}}(\mu_{\{1,\ldots,k\}})= (k+1)!\;$$
    where $\mu_{\{1,2,\ldots , k\}}$ is as defined in \cref{defn deformed affine frobenius realization}.
    
    \noindent We proceed by induction on $k$. The statement is true for the case $k=1$, since $\dd_1(x_1-\zeta x_2)=2$. For the general case, observe that    
    $$\mu_{\{1,\ldots,k\}} = \mu_{\{1,\ldots,k-1\}}(x_1-\zeta^kx_{k+1})(x_2-\zeta^{k-1}x_{k+1})\ldots (x_k -\zeta x_{k+1}).$$
    Since $(x_1-\zeta^kx_{k+1})(x_2-\zeta^{k-1}x_{k+1})\ldots (x_k -\zeta x_{k+1})$ is invariant for $W_{\{1,\ldots, k-1\}}$, we have
     \begin{align*}
     \dd_1\dd_2\ldots &\dd_k \dd_{\{1,2,\ldots , k-1\}}(\mu_{\{1,\ldots,k\}}) \\
     &= \dd_1 \dd_2 \ldots \dd_k\left(
     (x_1-\zeta^kx_{k+1})(x_2-\zeta^{k-1}x_{k+1})\ldots (x_k -\zeta x_{k+1})
          \dd_{\{1,2,\ldots , k-1\}}(\mu_{\{1,2,\ldots , k-1\}})
     \right)\\
     &=\dd_1\dd_2\ldots \dd_k\left((x_1-\zeta^kx_{k+1})(x_2-\zeta^{k-1}x_{k+1})\ldots (x_k -\zeta x_{k+1})\right)k!,
     \end{align*}
     so it suffices to show that
     \begin{equation} 
     \dd_1\dd_2\ldots \dd_k\left((x_1-\zeta^kx_{k+1})(x_2-\zeta^{k-1}x_{k+1})\ldots (x_k -\zeta x_{k+1})\right)=k+1.\label{eq 1 lemma-demazure standard rex}
     \end{equation}
    Observe that for any $s \in S$ $$\dd_s(f\al_s)= \frac{f\al_s+s(f)\al_s}{\al_s}=f+s(f).$$
    Taking $s=s_k$, we may simplify the LHS of \cref{eq 1 lemma-demazure standard rex} to
    \begin{align}
     \dd_1\dd_2\ldots &\dd_{k-1}\left(
     \begin{array}{c}
     (x_1-\zeta^kx_{k+1})(x_2-\zeta^{k-1}x_{k+1})\ldots (x_{k-1} -\zeta^2 x_{k+1})\vspace{5 pt}\\
     +\quad s_k((x_1-\zeta^k x_{k+1})(x_2-\zeta^{k-1}x_{k+1})\ldots (x_{k-1} -\zeta^2 x_{k+1}))
     \end{array}
     \right)  
     % \\
     % &=\notag \dd_1\dd_2\ldots \dd_{k-1}(x_1x_2\ldots x_{k-1})\\
     % &\;\;+\dd_1\dd_2\ldots \dd_{k-1}\left((x_1-\zeta^{k-1} x_{k})(x_2-\zeta^{k-2}x_{k})\ldots (x_{k-1} -\zeta x_{k})\right).
     \label{eq 2 lemma-demazure standard rex}
    \end{align}
    For the first term, using that $x_{k+1}$ is invariant under $s_i$ for $1 \le i \le k-1$, we get
    \begin{align*}\dd_1\dd_2\ldots \dd_{k-1}&((x_1-\zeta^kx_{k+1})(x_2-\zeta^{k-1}x_{k+1})\ldots (x_{k-1} -\zeta^2 x_{k+1}))\\
    &=\dd_1\dd_2\ldots \dd_{k-1}(x_1x_2\ldots x_{k-1}) + x_{k+1}(\text{something of negative degree})\\
    &=1+0=1.
    \end{align*}
    For the second term, by the inductive hypothesis,
    $$\dd_1\dd_2\ldots \dd_{k-1}\left((x_1-\zeta^{k-1} x_{k})(x_2-\zeta^{k-2}x_{k})\ldots (x_{k-1} -\zeta x_{k})\right)=k.$$
    So the RHS of \cref{eq 2 lemma-demazure standard rex} further simplifies to $1+k$ as desired.
\end{proof}

% For $I$ a finitary subset of $S$, the above result, along with  Lemma \ref{lemma deformed nilCoxeter}, tells us how to precisely relate $\partial_I$ to $\partial_{s_{i_1}} \ldots \partial_{s_{i_m}}$ for a reduced expression $(s_{i_1},s_{i_2},\ldots ,s_{i_m})$ of $w_I$. 
% We discuss this below.

\begin{defn}\label{defn-standard rex}
We define a \emph{standard rex}  for an arbitrary $I \subset S$ as follows.

When $I = \{1, 2, \ldots, k\}$ we define $\un{w}_I^{\std}$ as in the previous lemma, i.e.
\begin{equation}  \un{w}_{\{1, 2, \ldots,k\}}^{\std}:=(s_1,s_2,\ldots ,s_k)\bullet \un{w}^{\std}_{\{1,2, \ldots,k-1\}} \end{equation}
    where $\bullet$ is used for concatenation. Obviously $\un{w}_{\{1\}}^{\std}:=(s_1)$, the unique reduced expression.

When $I = \sigma^{l}(\{1, \ldots, k\})$, let $\un{w}_I^{\std} := \sigma^l(\un{w}_{\{1, \ldots, k\}}^{\std})$.

For $I$ disconnected,  let $I_1, \ldots, I_k$ denote the connected components of $I$ (in any order). Then define
\begin{equation} \un{w}_I^{\std} := \un{w}_{I_1}^{\std}\bullet \un{w}_{I_2}^{\std} \bullet \ldots \bullet \un{w}_{I_k}^{\std}. \end{equation}

\end{defn}

\begin{remark}
When $I$ is disconnected, the chosen order on components does change the expression $\un{w}_I^{\std}$ but not the operator $\dd_{\un{w}_I^{\std}}$, since $\dd_i$ and $\dd_j$ commute whenever $i$ and $j$ belong to different components.

\end{remark}

    Now we explore other reduced expressions.
    
\begin{defn} \label{defn-rex weight} For a finitary $I\subset S$, we define a function  $\mc{L}_I$ called the \emph{rex weight} on the set of reduced expressions of $w_I$. 
    First suppose $I=\set{1,2, \ldots, k}$ for some $k<n$.
    Let $\mathcal{L}_I'$ be the function on the collection of reduced expressions for $w_I$ such that 
    $$\mathcal{L}_I'((s_{i_1},s_{i_2},\ldots, s_{i_m}))=\sum_{j=1}^{m}i_j,$$
     and let $\mathcal{L}_I$ be the function such that 
    $$\LL_I((s_{i_1},s_{i_2},\ldots, s_{i_m}))= \LL_I'((s_{i_1},s_{i_2},\ldots, s_{i_m}))-\LL_I'(\un{w}_I^{\std}).$$
    In other words, $\LL_I$ counts the number of times we need to apply the the braid move $(s_i,s_{i+1},s_i)=(s_{i+1},s_i,s_{i+1})$ (with signs, i.e., +1 when going from $(s_i,s_{i+1},s_i)$ to $(s_{i+1},s_i,s_{i+1})$, and -1 the other way around) to get from the standard rex to the given reduced expression.

    For an arbitrary finitary $I \subset S$, we then define $\LL_I$ as follows.
    If $I$ is such that $\sigma^l(I)=\set{1,2, \ldots, k}$, we set $$\LL_I(\un{w}_I):=\LL_{\sigma^l(I)}(\sigma^l(\un{w}_I)).$$
    If $I$ is disconnected, with connected components $I_1, \ldots ,I_m$, we set
    $$\LL_I(\un{w}_I):=\sum_{i=1}^m\LL_{I_i}(\un{w}_{I_i})$$
    where $\un{w}_{I_j}$ are chosen such that
    $\un{w}_{I_1}\bullet \ldots \bullet \un{w}_{I_m}$ is a reduced expression of $w_I$ obtained from $\un{w}_I$ by only using braid moves
    of the form $(s_i,s_j)\sim(s_j,s_i)$.
\end{defn}

\begin{defn}\label{defn-rex weight2}
    For a subset $J$ of $I$ and a reduced expression $\un{w}_I^J$ of $w_Iw_J^{-1}$, we define the \emph{rex weight} $\LL_I^J(\un{w}_I^J)$ of $\un{w}_I^J$ to be $\LL_I^J(\un{w}_I^J):=\LL_I(\un{w}_I^J \bullet \un{w}_J^{\std})$.
\end{defn}
    
    Now \cref{lemma demazure operators computation deformed affine} (3) and Lemma \ref{lemma- demazure standard rex} quickly imply the following.
    
\begin{corollary}\label{corollary-normalization}
    For any $I\subset S$ finitary, and any reduced expression $(s_{i_1},s_{i_2},\ldots, s_{i_m})$ for $w_I$, we have
    $$\partial_I=\zeta^{-\LL_I((s_{i_1},s_{i_2},\ldots, s_{i_m})}\partial_{s_{i_1}}\ldots \partial_{s_{i_m}}. $$
    In particular, we have $\dd_I=\dd_{\un{w}_I^{\std}}.$

    Hence for $J\subset I$, and a reduced expression $\un{w}_I^J$ of $w_Iw_J^{-1}$, we have that
    $$\dd_I^J=\zeta^{-\LL_I^J(\un{w}_I^J)}\dd_{\un{w}_I^J}$$
\end{corollary}
\begin{remark} \label{rmk:wontworkforIJ}
    In the above setup, if $J\neq \emptyset$, it is not always true that there is a nice reduced expression $\un{w}_I^J$ of $w_Iw_J^{-1}$ so that $\dd_I^J=\dd_{\un{w}_I^J}$.
    For example, consider $J={2} \subset I=\set{1,2}$ for any $n>2$ (so that these are finitary). 
    Then $\dd_J=\dd_{s_2}$ and $\dd_I=\dd_{(s_1,s_2,s_1)}=\dd_{s_1}\dd_{s_2}\dd_{s_1}$. 
    But  $w_Iw_J^{-1}=s_2s_1$ has a unique reduced expression $(s_2,s_1)$ in this case, and $\dd_I^J=\zeta^{-1}\dd_{s_2}\dd_{s_1}$.
\end{remark}



%==========================================================================
\subsection{Calculations needed for geometric Satake}
\label{subsec-further frob calculations}
%==========================================================================

This section consists of various technical identities involving the Frealization structure on the deformed affine realization. We do not attempt to motivate each computation here; they will each be used in the verification of the diagrammatic Quantum Geometric Satake functor in \cref{sec well-defined-ness of the functor}. After \cref{subsubsec quantum numbers} and \cref{subsubsec-notation in further calculations.} the reader is welcome to skip this section.


\subsubsection{Quantum Numbers}\label{subsubsec quantum numbers}
For $n \in \Z_{\geq 1}$, we shall use the notation $[n]_q$ to refer to the \emph{n-th quantum number} explicitly given by 
$$[n]_q=\frac{q^n-q^{-n}}{q-q^{-1}}=q^{n-1}+q^{n-3}+\ldots q^{3-n}+q^{1-n} \in \Z[q,q^{-1}].$$
We also set $[0]_q=0$.
Thus $[1]_q=1, [2]_q=q+q^{-1}$. 
We also have explicit formulas for product of quantum numbers, for example 
$$[2]_q[n]_q=[n+1]_q+[n-1]_q.$$
More generally, for $m \geq n$, we have 
$$[m]_q[n]_q=[n]_q[m]_q=[m+n-1]_q+[n+m-3]_q+\ldots +[m-n+1]_q.$$
We often drop the subscript $q$ for convenience.
For $k \in \Z_{\geq 0}$, we define
$[k]!:= [k][k-1]\ldots [1]$.
We then define the \emph{quantum binomial coefficients} $\qbinom{n}{k}$ to be
$$\qbinom{n}{k}:= \frac{[n]!}{[n-k]! [k]!}.$$

% For the calculations below we introduce a formal variable $q$ satisfying 
% \begin{equation} \label{zetaq} \zeta^n = q^{-2}. \end{equation} We introduce $q$ so that we can view certain scalars as quantum numbers, and thus connect our calculations to calculations arising for quantum groups. Every time we use the variable $q$ it will be through one of the following two equalities:
% \begin{equation} \sum_{i=0}^{k} \zeta^{-in} = q^k [k+1], \end{equation}
% \begin{equation} \label{thisoneisworse} \zeta^{-\frac{n(n-1)}{2}} = q^{n-1}. \end{equation}
% There is no need to introduce $q$, as all scalars we use live in the ring $\Z[\zeta,\zeta^{-1}]$. Indeed, we only use the subring of $\Z[q,q^{-1}]$ generated by $q^{\pm 2}$ and $q^{n-1}$, all interpreted within $\Z[\zeta,\zeta^{-1}]$ via \eqref{zetaq} or \eqref{thisoneisworse}. Not all quantum numbers in $q$ live in this subring, but $[n]_q$ does.

For the calculations below we introduce a formal variable $q$ satisfying 
\begin{equation} \label{zetaq} \zeta^n = q^{-2}. \end{equation} We introduce $q$ so that we can view certain scalars as quantum numbers, and thus connect our calculations to calculations arising for quantum groups. 
More precisely, we work over 
\begin{equation}\label{eq coefficient ring with q and zeta}
    A_{\Z}:= \begin{cases}
\Z[q^{\pm 1}, \zeta^{\pm 1}]/(q^{-2}-\zeta^n) &\text{if $n$ odd.}\\
\Z[q^{\pm 1}, \zeta^{\pm 1}]/(q^{-1}-\zeta^{n/2}) &\text{if $n$ even.}
    \end{cases}
\end{equation}
Every time we use the variable $q$ it will be through one of the following two equalities:
\begin{equation} \sum_{i=0}^{k} \zeta^{-in} = q^k [k+1], \end{equation}
\begin{equation} \label{thisoneisworse} \zeta^{-\frac{n(n-1)}{2}} = q^{n-1}. \end{equation}
There is no need to introduce $q$, as all scalars we use live in the ring $\Z[\zeta,\zeta^{-1}]$. Indeed, we only use the subring of $\Z[q,q^{-1}]$ generated by $q^{\pm 2}$ and $q^{n-1}$, all interpreted within $\Z[\zeta,\zeta^{-1}]$ via \eqref{zetaq} or \eqref{thisoneisworse}. Not all quantum numbers in $q$ live in this subring, but $[n]_q$ does.


%To avoid confusion let us pin down $q$ formally. When $n$ is even, set $q = \zeta^{-n/2}$. When $n$ is odd, let $q$ be a formal square root of $\zeta^{-n}$. Then \eqref{thisoneisworse} holds, and it does not matter which square root is chosen when $n$ is odd, as replacing $q$ with $-q$ will not affect \eqref{thisoneisworse}.

\subsubsection{Some notation} \label{subsubsec-notation in further calculations.}
%We introduce some notation used in the rest of this section.
For $k,l,m \in S$, we continue using $\hh{k}$ to denote the (maximal finitary) subset $S \setminus \set{s} \subset S$, and $\hh{klm}$ to denote the subset $S \setminus \set{k,l,m}$, etc. 
We then have the corresponding Frobenius trace maps $\dd_{\hh{k}}$, $\dd_{\hh{klm}}$ and the product-coproduct elements $\mu_{\hh{k}}$, $\mu_{\hh{klm}}$.

For finitary $I \subset S$, $i,j \in I$ with $i\neq j$, we set
\begin{equation} \label{eq mu I I minus j I minus i}
  \mu_{I}^{I\setminus \set{i}, I\setminus \set{j}}:=  \frac{\mu_I\mu_{I\setminus \set{i,j}}}{\mu_{I\setminus \set{i}}\mu_{I\setminus \set{j}}}=\frac{\mu_I^{I\setminus \set{i}}}{\mu_{I\setminus \set{j}}^{I\setminus \set{i,j}}}.
\end{equation}
By the inclusion-exclusion principle, this is the product of all the positive roots in $\Phi^+_I$ which are not in $\Phi^+_{I \setminus \set{i}}$ or $\Phi^+_{I \setminus \set{j}}$.
For example, 
\begin{equation} \label{mu01k} \mu_{\hh{1}}^{\hh{01},\hh{1k}}= \frac{\mu_{\hh{1}}^{\hh{01}}}{\mu_{\hh{1k}}^{\hh{01k}}}= \frac{\prod_{2\leq i \leq n}(x_i - \zeta^{n-i}x_1)}{\prod_{k+1\leq i \leq n}(x_i - \zeta^{n-i}x_1)}=\prod_{2\leq i \leq k}(x_i - \zeta^{n+1-i}x_1).
\end{equation}
% This ratio is important when describing diagrammatics for $\SBSBim$ (\cref{subsec-diagrammatic singular bott-samelson bimodules}) and appears in the computations.

We also use Sweedler notation: having a $\Delta^J_{I}{}_{(1)}$ and $\Delta^J_{I}{}_{(2)}$ in the same expression indicates a sum over a dual basis $\set{e_i}, \set{f_i}$ for the extension $R^I \subset R^J$, i.e., we substitute $e_i$ for $\Delta^J_{I}{}_{(1)}$ and $f_i$ for $\Delta^J_{I}{}_{(2)}$ and take a sum over $i$.
For example, $\Delta^J_{I}{}_{(1)}\otimes \Delta^J_{I}{}_{(2)} = \sum_i e_i \ot f_i$ is by definition the image of $1$ under the coproduct map  $\Delta_I^J$ (see \cref{subsec- cubes of frob}).


\begin{remark}
    Using Sweedler notation, the expression in \cref{eq mu I I minus j I minus i} is also equal to 
    \begin{equation}  \mu_{I}^{I\setminus \set{i}, I\setminus \set{j}} =  \dd_{I\setminus \set{j}}^{I\setminus \set{i,j}}(\Delta^{I\setminus \set{i}}_{I,(1)})\Delta^{I\setminus \set{i}}_{I,(2)} = \dd_{I\setminus \set{i}}^{I\setminus \set{i,j}}(\Delta^{I\setminus \set{j}}_{I,(1)})\Delta^{I\setminus \set{j}}_{I,(2)}. \end{equation} This is a consequence of \cref{lemma deformed affine frobenius satisfies extra conditions}.
\end{remark}






%==========================================================================
\subsubsection{Dual bases}
\label{subsec-dual basis}

In this section, we construct explicit dual bases for certain extensions, which are necessary for further computations.

\begin{lemma} Given dual bases for $R^I \subset R^J$, we get dual bases for $R^{\sigma(I)} \subset R^{\sigma(J)}$ by applying $\sigma$ to the bases. \end{lemma}

\begin{proof} Immediate from \cref{rmk symmetry for rotation operator}. \end{proof}


\begin{lemma}\label{lemma dual basis left ext} The sets 
   $$ \{(-1)^{k-i}h_i(x_1)\}_{0\leq i \leq k},\;\{e_{k-i}(\zeta x_2, \zeta^2 x_3, \ldots, \zeta^{k}x_{k+1}) \}_{0\leq i \leq k}$$ form dual bases for the Frobenius extensions $R^{\{1,2, \ldots, k\}} \hookrightarrow R^{\{2, \ldots, k\}}$ for $k<n$. 
 
\end{lemma}
\begin{proof}
    We want to show:
    \begin{equation}\label{eq 1 dual basis left ext} 
        \dd_{\set{1,2, \ldots,k}}^{\set{2,3, \ldots, k}}\left((-1)^{k-i}h_i(x_1)e_{k-j}(\zeta x_2, \zeta^2 x_3, \ldots, \zeta^{k}x_{k+1})\right)=\delta_{i,j}.
    \end{equation}
    By \cref{corollary-normalization}, we have that 
    \begin{equation} \label{eq 2 dual basis left ext}
        \dd_{\set{1,2, \ldots,k}}^{\set{2,3, \ldots, k}}=\zeta^{-l}\dd_k \dd_{k-1} \ldots \dd_1
    \end{equation}
    where $l$ is the rex weight of $(s_k, s_{k-1}, \ldots s_1)$.
    Explicitly, from \cref{defn-rex weight}, 
    $$l=\LL_{\set{1,2,\ldots k}}((s_k,s_{k-1},\ldots, s_1)\bullet (s_2, s_3, \ldots, s_k, s_2, \ldots, s_{k-1}, \ldots, s_2, s_3, s_2) ).$$
    Comparing with the standard rex $\un{w}_{\set{1,2,\ldots,k}}=(s_1,s_2, \ldots s_k, s_1, \ldots ,s_{k-1}, \ldots, s_1, s_2, s_1)$, we see that
    $l=k(k-1)/2$.
    Hence \cref{eq 1 dual basis left ext} becomes
    \begin{equation}\label{eq 3 dual basis left ext}
        \zeta^{-k(k-1)/2}\dd_k \dd_{k-1} \ldots \dd_1\left((-1)^{k-i}h_i(x_1)e_{k-j}(\zeta x_2, \zeta^2 x_3, \ldots, \zeta^{k}x_{k+1})\right)=\delta_{i,j}.
    \end{equation}
    Now let's proceed to prove \cref{eq 3 dual basis left ext}. We use induction on $k$ for $1\leq k \leq n-1$. The base case is easy and left as an exercise to the reader. So now consider $k>1$.\\
    We first consider the case $i=0$. 
    In this case, if $j>0$, due to degree reasons, we get 
    $$\dd_k \dd_{k-1} \ldots \dd_1\left((-1)^{k}e_{k-j}(\zeta x_2, \zeta^2 x_3, \ldots, \zeta^{k}x_{k+1})\right)=0.$$
    When $j=0$, we have
    \begin{align}
        \nonumber&\dd_k \dd_{k-1} \ldots \dd_1\left((-1)^{k}e_{k}(\zeta x_2, \zeta^2 x_3, \ldots, \zeta^{k}x_{k+1})\right)\\
        &\nonumber= (-1)^k\dd_k \ldots \dd_1
        \left(
         \zeta x_2 e_{k-1}(\zeta^2x_3,\ldots, \zeta^kx_{k+1})
        +e_{k}(\zeta^2x_3, \ldots \zeta^k x_{k+1})        
        \right) \\
        &\nonumber=(-1)^k\dd_k\ldots\dd_2\left(\dd_1(\zeta x_2)\;e_{k-1}(\zeta^2x_3, \ldots, \zeta^kx_{k+1}) + 0\right)\\
        &=(-1)^{k-1}\dd_k\ldots \dd_2\left(\zeta^{k-1}e_{k-1}(\zeta x_3, \ldots, \zeta^{k-1} x_{k+1})\right)\label{eq 4 dual basis left ext}\\
        &=(-1)^{k-1}\zeta^{k-1}\dd_{k-1}\ldots \dd_1\left(e_{k-1}(\zeta x_2, \ldots, \zeta^{k-1} x_{k})\right) \label{eq 5 dual basis left ext} \\
        &=\zeta^{k-1} \zeta^{(k-1)(k-2)/2}=\zeta^{k(k-1)/2} \label{eq 6 dual basis left ext}
    \end{align}
    and we have that \cref{eq 3 dual basis left ext} holds. Note that in the above computation, we used symmetry with respect to the rotation operator $\sigma$, (see \cref{rmk symmetry for rotation operator}) to go from \eqref{eq 4 dual basis left ext} to \eqref{eq 5 dual basis left ext}, and induction hypothesis to go from \eqref{eq 5 dual basis left ext} to \eqref{eq 6 dual basis left ext}.

    Now suppose $i>0$. Then we have: 
    \begin{align}
        \nonumber(-1)^{k-i}&\dd_k\ldots \dd_2 \dd_1\left(x_1^ie_j(\zeta x_2, \ldots , \zeta^k x_{k+1}\right) \\
        \nonumber&= (-1)^{k-i}\dd_k \ldots \dd_2\dd_1 \left(x_1^i\zeta x_2 e_{j-1}(\zeta^2x_3, \ldots, \zeta^k x_{k+1})+x_1^ie_j(\zeta^2 x_3, \ldots , \zeta^k x_{k+1}) \right)\\
       \nonumber &= (-1)^{k-i}\dd_k \ldots \dd_2\left(x_1 \zeta x_2 \dd_1(x_1^{i-1})e_{j-1}(\zeta^2 x_3, \ldots , \zeta^k x_{k+1}) + \dd_1(x_1^i)e_j(\zeta^2 x_3, \ldots , \zeta^k x_{k+1})
        \right)\\
        \nonumber&= (-1)^{k-i} \dd_k \ldots \dd_2\left(
        \begin{array}{cc}
         &\displaystyle x_1 \zeta x_2 \sum_{l'=0}^{i-2} x_1^{i-2-l'}(\zeta x_2)^{l'}e_{j-1}(\zeta^2 x_3, \ldots \zeta^k x_{k+1})    \\
          + &\displaystyle \sum_{l=0}^{i-1}x_1^{i-1-l}(\zeta x_2)^le_j(\zeta^2x_3, \ldots, \zeta^k x_{k+1})   
        \end{array}\right)\\
        &= (-1)^{k-i}\left(
        \begin{array}{cc}
        &\displaystyle \sum_{l'=0}^{i-2} x_1^{i-1-l'}\dd_k \ldots \dd_2((\zeta x_2)^{l'+1}e_{j-1}(\zeta^2 x_3, \ldots \zeta^k x_{k+1}))    \\
          + &\displaystyle \sum_{l=0}^{i-1}x_1^{i-1-l}\dd_k \ldots \dd_2((\zeta x_2)^le_j(\zeta^2x_3, \ldots, \zeta^k x_{k+1}))
        \end{array}\right)
    \end{align}
    Since $i,j \leq k$, we have that $i-1, j-1 \leq k-1$. 
    An argument using rotational symmetry and induction hypothesis shows then that
    \begin{align*}\dd_k \ldots \dd_2((\zeta x_2)^{l'+1}e_{j-1}(\zeta^2 x_3, \ldots \zeta^k x_{k+1}))&=\dd_{k-1} \ldots \dd_1((\zeta x_1)^{l'+1}e_{j-1}(\zeta^2 x_2, \ldots, \zeta^k x_{k}))\\
    &=\zeta^{l'+j}\dd_{k-1} \ldots \dd_1((x_1)^{l'+1}e_{j-1}(\zeta x_2, \ldots ,\zeta^{k-1} x_{k}))\\
    &=\zeta^{k-1}(-1)^{k-l'-2}\zeta^{(k-1)(k-2)/2}\delta_{j-1,k-l'-2}\\
    &=\zeta^{k(k-1)/2}(-1)^{j-1}\delta_{j-1,k-l'-2},
    \end{align*}
    and
    \begin{align*}
     \dd_k \ldots \dd_2((\zeta x_2)^le_j(\zeta^2x_3, \ldots, \zeta^k x_{k+1}))&=\dd_{k-1} \ldots \dd_1((\zeta x_1)^le_j(\zeta^2x_2, \ldots, \zeta^k x_{k}))  \\
     &= \zeta^{j+l} \dd_{k-1} \ldots \dd_1(( x_1)^le_j(\zeta x_2, \ldots, \zeta^{k-1} x_{k}))\\
     &=\zeta^{j+l}(-1)^{k-1-l}\zeta^{(k-1)(k-2)/2}\delta_{j, k-1-l}\\
     &=\zeta^{k(k-1)/2}(-1)^{j}\delta_{j,k-1-l}.
    \end{align*}
Hence RHS of \eqref{eq 6 dual basis left ext} simplifies to
\begin{align*}
    &(-1)^{k-i}\left(
    \begin{array}{c}
    \displaystyle \sum_{l'=0}^{i-2} x_1^{i-1-l'}\zeta^{k(k-1)/2}(-1)^{j-1}\delta_{j-1,k-l'-2}    \\
    + \displaystyle \sum_{l=0}^{i-1}x_1^{i-1-l}\zeta^{k(k-1)/2}(-1)^{j}\delta_{j,k-1-l}
     \end{array}   \right)\\
    &=(-1)^{k-i}\zeta^{k(k-1)/2}(-1)^{j} \left(
    \displaystyle \sum_{l=0}^{i-1}x_1^{i-1-l}\delta_{j,k-1-l}
    -\displaystyle \sum_{l'=0}^{i-2} x_1^{i-1-l'}\delta_{j-1,k-l'-2}
    \right)\\
    &=(-1)^{k-i-j}\zeta^{k(k-1)/2}\delta_{j,k-i}=\zeta^{k(k-1)/2}.
\end{align*}
This implies \eqref{eq 3 dual basis left ext}.
\end{proof}

\begin{remark}\label{rmk explicit dual basis permutation realization}
    \cref{lemma dual basis left ext} with $\zeta$ specialized to $1$ provides explicit dual basis for the corresponding extension for the permutation realization of $S_n$ over $\Z$. 
\end{remark}

% For checking the well-definedness of the functor $\GS_{\zeta}$ in \cref{sec well-defined-ness of the functor}, we would also need an explicit dual basis for the extension $R^{\{1,\ldots, k-1\}} \hookrightarrow R^{\{1,\ldots k\}}$, hence the following corollary.

\begin{corollary}\label{corollary dual basis right ext} The sets
     $$ \{(-1)^{i}\zeta^{k(k+1)/2}h_i(x_{k+1})\}_{0\leq i \leq k},\;\{e_{k-i}(\zeta^{-1} x_k, \zeta^{-2} x_{k-1}, \ldots, \zeta^{-k}x_1) \}$$ form dual bases for the Frobenius extensions $R^{\{1,2, \ldots, k\}} \hookrightarrow R^{\{1,2, \ldots, k-1\}}$ for $k<n$.
     % A similar statement holds for any of the extensions $R^{\set{l+1, l+2, \ldots, l+k}}\hookrightarrow R^{\set{l+1, l+2, \ldots l+k-1}}$ for any $l\in \Z$, with indices shifted accordingly (see \cref{rmk symmetry for rotation operator}).
\end{corollary}
\begin{proof}
    One can do a similar explicit proof as in the proof of \cref{lemma dual basis left ext} to show this. 
    Alternatively, one may use \cref{rmk dual basis from permutation realization} and \cref{rmk explicit dual basis permutation realization} to see this as a consequence of \cref{lemma dual basis left ext}. 
    One may also see this as a consequence of \cref{lemma dual basis left ext} using symmetry of our Frobenius realization under the rotation \eqref{eq rotational symmetry deformed affine} and reflection \eqref{eq:reflection symmetry for reflection operator}.
\end{proof}




\subsubsection{Calculations}\label{subsubsec calculations}
% We now proceed to the computations involving the Frobenius realization structure on the deformed affine realization. 

\begin{lemma}\label{lemma demazure adjacent} Let $0\leq k<n$. Then we have the following equalities (whenever all subsets appearing are finitary):  
    \begin{enumerate}
        \item[(a)]   \hfill$ \dd_{\{1,2, \ldots ,k\}}^{\{1,2, \ldots ,k-1\}}\left((x_1-\zeta^{-(n-k)}x_{k+1})\ldots (x_{k}-\zeta^{-(n-1)}x_{k+1})\right)=q^{k}[k+1]_q$ \hfill \null
        \item[(b)] \hfill $\dd^{\hh{0,-1}}_{\hh{0}}(\mu^{\hh{0,-1}}_{\hh{-1}}) =(-1)^{n-1}[n]_q$. \hfill \null
        \item[(c)]  \hfill $\dd_{\{1,2, \ldots ,k-1\}}^{\{2, \ldots ,k-1\}}\left((x_1-\zeta^{-(n-k+1)}x_k)\ldots (x_1-\zeta^{-(n-1)}x_2)\right)=q^{k-1}[k]_q.$ \hfill \null\label{lemma demazure for bigon}
        \item[(d)] \hfill $ \dd_{\hh{0}}^{\hh{0,1}}\left(\mu^{\hh{0,1}}_{\hh{1}} \right)= (-1)^{n-1}[n]_q$. \hfill \null
    \end{enumerate}
\end{lemma}
\begin{proof}   
  To prove part (a), observe that 
  \begin{align*}(x_1-\zeta^{-(n-k)}x_{k+1})\cdots &(x_{k}-\zeta^{-(n-1)}x_{k+1})\\
  &= (\zeta^{-(n-k)})\cdots (\zeta^{-(n-1)})(\zeta^{n-k}x_1-x_{k+1})\cdots (\zeta^{n-1}x_k-x_{k+1}) \\
  &=\zeta^{-kn+\frac{k(k+1)}{2}}\sum_{0\leq i \leq k}(-1)^ih_i(x_{k+1})e_{k-i}(\zeta^{n-1}x_k, \ldots, \zeta^{n-k}x_1)\\
  &=q^{2k}\sum_{0 \leq i \leq k}(-1)^i\zeta^{\frac{k(k+1)}{2}}h_{i}(x_{k+1})\zeta^{n(k-i)}e_{k-i}(\zeta^{-1}x_k, \ldots, \zeta^{-k}x_1).
  \end{align*}
  Applying $\dd_{\{1,2, \ldots ,k\}}^{\{1,2, \ldots ,k-1\}}$ to both sides, using \cref{corollary dual basis right ext}, we see that
  \begin{align*}
    \dd_{\{1,2, \ldots ,k\}}^{\{1,2, \ldots ,k-1\}}\left((x_1-\zeta^{-(n-k)}x_{k+1})\ldots (x_{k}-\zeta^{-(n-1)}x_{k+1})\right)  
    =q^{2k}\sum_{0 \leq i \leq k}q^{-2(k-i)}=q^{k}[k+1]_q.
  \end{align*}
  
  Now we show that (a) implies (b).
    Observe that from \cref{defn deformed affine frobenius realization} and \cref{lemma deformed affine frobenius realization is frobenius}, we have that
    $$\mu_{\{0,1, \ldots ,n-2\}}^{\{1,2, \ldots ,n-2\}}=(x_n-\zeta x_1)\ldots (x_n-\zeta^{n-1} x_{n-1}).$$
    So (c) reduces to showing
    $$\dd_{\{1,2, \ldots ,n-1\}}^{\{1,2, \ldots ,n-2\}}((x_n-\zeta x_1)\ldots (x_n-\zeta^{n-1} x_{n-1}))=(-1)^{n-1}[n]_q.$$
    We can rewrite the L.H.S above as:
    \begin{align}       
   LHS &= (-\zeta)(-\zeta^2)\ldots (-\zeta^{n-1})\dd_{\{1,2, \ldots ,n-1\}}^{\{1,2, \ldots ,n-2\}}((x_1-\zeta^{-1}x_n)\ldots (x_{n-1}-\zeta^{-n+1}x_n)) \nonumber\\
   &= (-1)^{n-1}\zeta^{\frac{n(n-1)}{2}}\dd_{\{1,2, \ldots ,n-1\}}^{\{1,2, \ldots ,n-2\}}((x_1-\zeta^{-1}x_n)\ldots (x_{n-1}-\zeta^{-n+1}x_n))\nonumber\\
   &= (-1)^{n-1}q^{-(n-1)}\dd_{\{1,2, \ldots ,n-1\}}^{\{1,2, \ldots ,n-2\}}((x_1-\zeta^{-1}x_n)\ldots (x_{n-1}-\zeta^{-n+1}x_n)).\nonumber
   \end{align}
   Hence (b) follows from (a), taking $k$ to be $n-1$.

   Now we prove (c).
   LHS may be written as
    \begin{align*}
    &\dd_{\{1,2, \ldots ,k-1\}}^{\{2, \ldots ,k-1\}}\left((x_1-\zeta^{-(n-k+1)}x_k)\ldots (x_1-\zeta^{-(n-1)}x_2)\right) \\
    &\quad= \dd_{\{1,2, \ldots ,k-1\}}^{\{2, \ldots ,k-1\}}\left(\sum_{0 \leq i \leq k-1}h_i(x_1)e_{k-1-i}(-\zeta^{-n+1}x_2, -\zeta^{-n+2} x_3, \ldots, -\zeta^{-n+(k-1)}x_k)\right)\\
    &\quad=\dd_{\{1,2, \ldots ,k-1\}}^{\{2, \ldots ,k-1\}}\left(\sum_{0 \leq i \leq k-1}(-1)^{k-1-i}\zeta^{-n(k-1-i)}h_i(x_1)e_{k-1-i}(\zeta x_2, \zeta^{2} x_3, \ldots, \zeta^{-(k-1)}x_k)\right)\\
    &\quad = \sum_{0 \leq i \leq k-1} \zeta^{-n(k-1-i)} = \sum_{0 \leq i \leq k-1}q^{2(k-1-i)}=q^{(k-1)}[k]_q.
    \end{align*}
    where we used \cref{lemma dual basis left ext} for the first equality in the last line above.
    For (d), we have from \cref{defn deformed affine frobenius realization} that
    $$\mu^{\hh{0,1}}_{\hh{1}}= (x_2-\zeta^{n-1}x_1)(x_3-\zeta^{n-2}x_1)\ldots (x_n-\zeta x_1)=(-1)^{n-1}\zeta^{\frac{n(n-1)}{2}}\cdot(x_1-\zeta^{-1}x_{n})\ldots (x_1-\zeta^{-(n-1)}x_2).$$
    Noting that $\zeta^{n(n-1)/2}= q^{-(n-1)},$ (d) then follows from (c) when $k=n.$  
\end{proof}




% \comm{
% \begin{lemma}
%     $\dd_{S\bs \{0\}}^{S\bs \{0,k\}}(\mu_{S\bs \{k\}}^{S\bs \{0,k\}})= (-1)^{k(n-k})\qbinom{n}{k}_q$ (for $SL_n$.)
% \end{lemma}
% \begin{proof}
%     \BE{to fix later...} I didn't really write this down- but with Stephen Begalow's presentation with smaller number of relations, this will follow once we construct the diagrammatic Satake functor over $\mathbb{Q}(q,\zeta)$.
% \end{proof}}


\begin{corollary}\label{corollary- bigon bursting}
For $k \not= 0,1$, we have:
    $$\dd_{\hh{0,k}}^{\hh{0,1,k}}\left(\mu_{\hh{1}}^{\hh{0,1},\hh{1,k}}\right)=(-1)^{k-1}\zeta^{-k(k-1)/2}q^{-(k-1)}[k]_q.$$
    \begin{proof}
        Rewriting the LHS, we get
        \begin{align}
        \dd_{\hh{0,k}}^{\hh{0,1,k}}\left(\mu_{\hh{1}}^{\hh{0,1},\hh{1,k}}\right)&=\dd_{\hh{0,k}}^{\hh{0,1,k}}\left( (x_2-\zeta^{n-1}x_1)(x_3-\zeta^{n-2}x_1)\ldots (x_k-\zeta^{n+1-k}x_1)\right). \label{eq 5.7.1}
        \end{align}
        We also have that 
        \[\dd_{\hh{0,k}}^{\hh{0,1,k}}=\dd_{\set{1,\ldots, k-1}\sqcup\set{k+1,\ldots, n-1}}^{\set{2,\ldots, k-1}\sqcup \set{k+1, \ldots, n-1}}=\dd_{\set{1,\ldots, k-1}}^{\set{2,\ldots, k-1}}.\]
       Hence we may simplify \eqref{eq 5.7.1} further:
       \begin{align}
        \dd_{\hh{0,k}}^{\hh{0,1,k}}\left(\mu_{\hh{1}}^{\hh{0,1},\hh{1,k}}\right)&=\dd_{\set{1,\ldots, k-1}}^{\set{2,\ldots, k-1}}\left( (x_2-\zeta^{n-1}x_1)(x_3-\zeta^{n-2}x_1)\ldots (x_k-\zeta^{n+1-k}x_1)\right) \nonumber\\
        &= (-1)^{k-1}\zeta^{(n-1)+(n-2)+\ldots + (n+1-k)}
        \\&\hspace{20pt}\dd_{\set{1,\ldots, k-1}}^{\set{2,\ldots, k-1}}\left((x_1-\zeta^{-(n-k+1)}x_k)\ldots (x_1-\zeta^{-(n-1)}x_2)\right)\nonumber\\
        &=(-1)^{k-1}\zeta^{(k-1)n-\frac{(k-1)k}{2}}q^{(k-1)}[k]_q \label{eq 5.7.2}\\
        &=(-1)^{k-1}(q^{-2})^{k-1}\zeta^{-\frac{(k-1)k}{2}}q^{(k-1)}[k]_q = (-1)^{k-1}\zeta^{-\frac{(k-1)k}{2}}q^{-(k-1)}[k]_q \label{eq 5.7.3}
        \end{align}
        where we used Lemma \ref{lemma demazure for bigon} for the equality in \eqref{eq 5.7.2} and the fact that $\zeta^n=q^{-2}$ for the first equality in \eqref{eq 5.7.3}. 
    \end{proof}
\end{corollary}

\begin{lemma}\label{lemma square flop 1 tensor 1}
    The following equality holds in $\Bimod{\hh{1,k+1}}{}{R^{\hh{1,k+1}}}\otimes_{\hh{1}}\Bimod{}{\hh{0,1}}{R^{\hh{0,1}}}$ for $2\leq k\leq n-1$.
    
    $$\ddm{1,2, k+1}{1,k+1}\left(\ddm{1,2,k+1}{2,k+1}(\Del{1,2}{2}{1})\Del{1,2}{1}{1}\right) \otimes \ddm{0,1,2}{0,1}\left(\ddm{0,1,2}{0,2}(\Del{1,2}{2}{2})\Del{1,2}{1}{2} \right) =[k-1]_q(-1)^k\zeta^{\frac{-(k-1)k}{2} -1}q^{-k}(1\otimes 1)$$
\end{lemma}
\begin{proof}
    Let us start by analyzing the LHS. 
    Observe that we need to make explicit choices for the dual basis for the extensions $R^{\hh{2}}\hookrightarrow R^{\hh{1,2}}$ and $R^{\hh{1}}\hookrightarrow R^{\hh{1,2}}$ to make the computations involving $\Del{1,2}{1}{.}$ and $\Del{1,2}{2}{.}$. 
    This is precisely the content of \cref{lemma dual basis left ext} and \cref{corollary dual basis right ext}; we choose to let $\Del{1,2}{2}{1}$ run over $\set{e_{n-1-i}(\zeta^{-1}x_1, \zeta^{-2}x_n, \ldots , \zeta^{-(n-1)}x_3)}$ and $\Del{1,2}{2}{2}$ run over $ \set{(-1)^i\zeta^{n(n-1)/2}h_i(x_2)}$.
    Similarly, we choose $\Del{1,2}{1}{1}$ run over $\set{(-1)^{n-1-j}h_{j}(x_2)}$ and $\Del{1,2}{1}{2}$ run over $\set{e_{n-1-j}(\zeta x_3, \zeta^2 x_4, \ldots \zeta^{n-1} x_1)}.$
    So we may write the LHS as:
    \begin{align}
    \sum_{0\leq i,j \leq n-1} \left( \begin{array}{c}
    \ddm{1,2, k+1}{1,k+1}\left(\ddm{1,2,k+1}{2,k+1}\left(e_{n-1-i}(\zeta^{-1}x_1, \zeta^{-2}x_n, \ldots , \zeta^{-(n-1)}x_3)\right)(-1)^{n-1-j}h_{j}(x_2)\right) \vspace{5 pt}\\ \otimes \ddm{0,1,2}{0,1}\left(\ddm{0,1,2}{0,2}((-1)^i\zeta^{n(n-1)/2}h_i(x_2))e_{n-1-j}(\zeta x_3, \zeta^2 x_4, \ldots \zeta^{n-1} x_1) \right) \label{eq 1 lemma square flop 1 tensor 1}
    \end{array}\right)
    \end{align}
    Note that the sum above should give a degree 0 term in $\Bimod{\hh{1,k+1}}{}{R^{\hh{1,k+1}}}\otimes_{\hh{1}}\Bimod{}{\hh{1,2}}{R^{\hh{1,2}}}$. This implies that we need only consider tuples $(i,j)$ where both the tensor factors are of degree $0$.
    This happens precisely when $i-j=0$, i.e., when $i=j$.
    So we now just analyze the terms in \eqref{eq 1 lemma square flop 1 tensor 1} for $i=j.$
    
    Now let us simplify the second tensor factor on the LHS. This factor is completely independent of $k$.
    We have that $\ddm{0,1,2}{0,2}=\dd^{\set{3, 4, \ldots , n-1}}_{\set{1}\sqcup \set{3,4, \ldots, n-1}}=\dd_1$, so we have that the $i=0$ term is 0. 
    Now we can simplify the second term in \eqref{eq 1 lemma square flop 1 tensor 1} to be
    \begin{align} \label{floobarl}
    &(-1)^i\zeta^{n(n-1)/2} \ddm{0,1,2}{0,1}\left( \dd_1(h_i(x_2))e_{n-1-i}(\zeta x_3, \zeta^2 x_4, \ldots , \zeta^{n-1}x_1) \right)\\ \nonumber
    &=(-1)^i\zeta^{n(n-1)/2} \ddm{0,1,2}{0,1}\left( (-\zeta^{-1}) h_{i-1}(x_2, \zeta^{-1}x_1)e_{n-1-i}(\zeta x_3, \zeta^2 x_4, \ldots , \zeta^{n-1}x_1)
    \right)\\ \nonumber
    &=(-1)^{i+1}\zeta^{\frac{n(n-1)}{2} -1} \ddm{0,1,2}{0,1}\left(
    \sum_{l=0}^{i-1}(\zeta^{-1}x_1)^{i-1-l}h_l(x_2)e_{n-1-i}(\zeta x_3, \zeta^2 x_4, \ldots , \zeta^{n-1}x_1)
    \right)\end{align}
    % &=(-1)^{i+1}\zeta^{\frac{n(n-1)}{2} -1} \ddm{0,1,2}{0,1}\left(
    % \sum_{l=0}^{i-1}(\zeta^{-1}x_1)^{i-1-l}h_l(x_2)
    % \left(\begin{array}{c}
    % e_{n-1-i}(\zeta x_3, \zeta^2 x_4, \ldots , \zeta^{n-2}x_0) \vspace{5pt}\\
    % + \zeta^{n-1}x_1e_{n-2-i}(\zeta x_3, \zeta^2 x_4, \ldots , \zeta^{n-2}x_0)
    % \end{array}\right)
    % \right).
    Now $x_1 \in W_{\hh{0,1}}$ so it can be pulled out of the operator $\dd^{\hh{0,1,2}}_{\hh{0,1}}$. If this happens then $\dd^{\hh{0,1,2}}_{\hh{0,1}}$ applied to the remainder would be zero, for degree reasons. Consequently we can ignore all terms above where $x_1$ appears as a factor, both those terms where $l \ne i-1$, and those terms within $e_{n-1-i}$ where $\zeta^{n-1} x_1$ is chosen as a factor. Consequently \eqref{floobarl} simplifies to
    % \begin{align*}
    % &(-1)^{i+1}\zeta^{\frac{n(n-1)}{2} -1} \ddm{0,1,2}{0,1}\left(
    % h_{i-1}(x_2)
    % \left(\begin{array}{c}
    % e_{n-1-i}(\zeta x_3, \zeta^2 x_4, \ldots , \zeta^{n-2}x_0) \vspace{5pt}\\
    % + \zeta^{n-1}x_1e_{n-2-i}(\zeta x_3, \zeta^2 x_4, \ldots , \zeta^{n-2}x_0)
    % \end{array}\right)
    % \right)\\
    % &=(-1)^{i+1}\zeta^{\frac{n(n-1)}{2}-1} (-1)^{n-2-(i-1)}(\delta_{i-1, n-2-(n-1-i)} +\zeta^{n-1}x_1\delta _{i-1, n-2 - (n-2-i)})\\
    % &=(-1)^{n}\zeta^{\frac{n(n-1)}{2}-1}.
    % \end{align*}
     \begin{align}
    &(-1)^{i+1}\zeta^{\frac{n(n-1)}{2} -1} \ddm{0,1,2}{0,1}\left(
    h_{i-1}(x_2)
    e_{n-1-i}(\zeta x_3, \zeta^2 x_4, \ldots , \zeta^{n-2}x_0)
    \right)\\ \nonumber
    &=(-1)^{i+1}\zeta^{\frac{n(n-1)}{2}-1} (-1)^{n-2-(i-1)}
    =(-1)^{n}\zeta^{\frac{n(n-1)}{2}-1}.
    \end{align}
    The first equality above used \cref{lemma dual basis left ext}.
    

    Hence \eqref{eq 1 lemma square flop 1 tensor 1} becomes
    \begin{align}
    \sum_{1\leq i \leq n-1} \left( \begin{array}{c}
    \ddm{1,2, k+1}{1,k+1}\left(\ddm{1,2,k+1}{2,k+1}\left(e_{n-1-i}(\zeta^{-1}x_1, \zeta^{-2}x_n, \ldots , \zeta^{-(n-1)}x_3)\right)(-1)^{n-1-i}h_{i}(x_2)\right) \vspace{5 pt}\\ \otimes (-1)^{n}\zeta^{\frac{n(n-1)}{2}-1} \label{eq 2 lemma square flop 1 tensor 1}
    \end{array}\right)
    \end{align}

    Now let us analyze the first tensor factor.
    We may write the Frobenius trace maps $\ddm{1,2, k+1}{1,k+1}$ and $\ddm{1,2,k+1}{2,k+1}$ explicitly using \cref{corollary-normalization} as
    \[\ddm{1,2,k+1}{2,k+1}=\dd^{\set{3, 4, \ldots , k} \sqcup \set{k+2, \ldots, n-1, n}}_{\set{3,4, \ldots k}\sqcup \set{k+2,\ldots,n,1}}=\dd_{k+2}\ldots \dd_n \dd_1, \]
    \[ \ddm{1,2,k+1}{1,k+1}=\dd^{\set{3, 4, \ldots , k} \sqcup \set{k+2, \ldots, n-1, n}}_{\set{2,3 \ldots k}\sqcup \set{k+2,\ldots,n-1, n}}=\zeta^{-\frac{(k-1)(k-2)}{2}}\dd_k \ldots \dd_3 \dd_2=\dd^{\set{3,4, \ldots, k}}_{\set{2, 3, \ldots, k}}.\]

    The first tensor factor in \eqref{eq 2 lemma square flop 1 tensor 1} (for a given $i\geq0$) then becomes:
    \begin{align}
    &\dd^{\set{3,4, \ldots, k}}_{\set{2, 3, \ldots, k}} \left(
    \dd_{k+2}\ldots \dd_n \dd_1 (e_{n-1-i}(\zeta^{-1}x_1, \zeta^{-2}x_n, \ldots \zeta^{-(n-1)}x_3))(-1)^{n-1-i}h_i(x_2)    
    \right). \label{eq 3 lemma square flop 1 tensor 1}
    \end{align}
    Every term in $e_{n-1-i}(\zeta^{-1}x_1, \zeta^{-2}x_n, \ldots \zeta^{-(n-1)}x_3))$ which does not include $\zeta^{-1} x_1$ is invariant under $s_1$ and thus killed by $\dd_1$. Thus
    \begin{align} \dd_1(e_{n-1-i}&(\zeta^{-1}x_1, \zeta^{-2}x_n, \ldots \zeta^{-(n-1)}x_3)) = e_{n-2-i}(\zeta^{-2}x_n, \ldots \zeta^{-(n-1)}x_3) \dd_1(\zeta^{-1} x_1) = \nonumber \\
    & \zeta^{-1} e_{n-2-i}(\zeta^{-2}x_n, \ldots \zeta^{-(n-1)}x_3). \end{align}
    Similarly, applying $\dd_n$ to $e_{n-2-i}(\zeta^{-2}x_n, \ldots \zeta^{-(n-1)}x_3)$ will kill any term not including $\zeta^{-2} x_n$, and so forth. We deduce that
    % \begin{align}
    %    \notag &\dd_{k+2}\ldots \dd_n \dd_1 (e_{n-1-i}(\zeta^{-1}x_1, \zeta^{-2}x_n, \ldots \zeta^{-(n-1)}x_3))\\
    %    \notag &=\dd_{k+2}\ldots \dd_n \dd_1 \left(\sum_{l=0}^{n-1-i} 
    %     \left(
    %     \begin{array}{l}
    %     e_{n-1-i-l}(\zeta^{-1}x_1, \zeta^{-2}x_n, \ldots \zeta^{-(n-k)}x_{k+2})\cdot \vspace{5 pt}\\
    %     \hspace{30 pt}e_{l}(\zeta^{-(n-k+1)}x_{k+1}, \ldots, \zeta^{-(n-2)}x_4, \zeta^{-(n-1)}x_3)
    %     \end{array}\right)
    %     \right)\\
    %    \notag &=\sum_{l=0}^{n-1-i}
    %     \left( \begin{array}{l}
    %     e_{l}(\zeta^{-(n-k+1)}x_{k+1}, \ldots, \zeta^{-(n-2)}x_4, \zeta^{-(n-1)}x_3) \cdot \vspace{5 pt}\\
    %     \hspace{30 pt}\dd_{k+2}\ldots \dd_n \dd_1 \left(e_{n-1-i-l}(\zeta^{-1}x_1, \zeta^{-2}x_n, \ldots \zeta^{-(n-k)}x_{k+2})\right)
    %     \end{array}\right)\\
    %    \notag &= \sum_{l=0}^{n-1-i} e_{l}(\zeta^{-(n-k+1)}x_{k+1}, \ldots, \zeta^{-(n-2)}x_4, \zeta^{-(n-1)}x_3) \zeta^{-\frac{(n-k)(n-k+1)}{2}}\delta_{n-1-i-l, n-k}  \hspace{5 pt}\text{(by \cref{corollary dual basis right ext})}\\
    %     &= \zeta^{-\frac{(n-k)(n-k+1)}{2}}e_{k-1-i}(\zeta^{-(n-k+1)}x_{k+1}, \ldots, \zeta^{-(n-2)}x_4, \zeta^{-(n-1)}x_3), \label{eq left tensor factor square flop}
    % \end{align}
       \begin{align}
       \nonumber &\dd_{k+2}\ldots \dd_n \dd_1 (e_{n-1-i}(\zeta^{-1}x_1, \zeta^{-2}x_n, \ldots \zeta^{-(n-1)}x_3))\\ \nonumber
       & = \zeta^{-1-2-\ldots -(n-k)} e_{k-1-i}(\zeta^{-(n-k+1)}x_{k+1}, \ldots, \zeta^{-(n-2)}x_4, \zeta^{-(n-1)}x_3) \\ 
        &= \zeta^{-\frac{(n-k)(n-k+1)}{2}}e_{k-1-i}(\zeta^{-(n-k+1)}x_{k+1}, \ldots, \zeta^{-(n-2)}x_4, \zeta^{-(n-1)}x_3). \label{eq left tensor factor square flop}
    \end{align}
    Thus \eqref{eq 3 lemma square flop 1 tensor 1} simplifies to
    \begin{align}
    \notag  &\dd^{\set{3,4, \ldots, k}}_{\set{2, 3, \ldots, k}} \left(
     \zeta^{-\frac{(n-k)(n-k+1)}{2}}e_{k-1-i}(\zeta^{-(n-k+1)}x_{k+1}, \ldots, \zeta^{-(n-2)}x_4, \zeta^{-(n-1)}x_3)
    (-1)^{n-1-i}h_i(x_2)  
    \right) \\
    \notag&=\zeta^{-\frac{(n-k)(n-k+1)}{2}} (-1)^{n-1-i} \dd^{\set{3,4, \ldots, k}}_{\set{2, 3, \ldots, k}} \left( \zeta^{-n(k-1-i)}e_{k-1-i}(\zeta x_3, \zeta^2 x_4, \ldots, \zeta^{k-1}x_{k+1})h_i(x_2)
    \right)\\
    &= \begin{cases}
    \begin{array}{cc}
    \zeta^{-\frac{(n-k)(n-k+1)}{2}}\zeta^{-n(k-1-i)} (-1)^{n-1-i} (-1)^{k-1-i}  &\text{if $0\leq i \leq k-1$} \\
    0  &\text{otherwise} 
    \end{array}
    \end{cases} \label{eq 4 lemma square flop 1 tensor 1}\\
    &=\begin{cases}
    \begin{array}{cc}
    \zeta^{-\frac{(n-k)(n-k+1)}{2}-n(k-1-i)} (-1)^{n-k}  &\text{if $0\leq i \leq k-1$}\\
    0  &\text{otherwise}
    \end{array}
    \end{cases}
    \label{eq 5 lemma square flop 1 tensor 1}  
    \end{align}
    where we used \cref{lemma dual basis left ext} for the equality in \eqref{eq 4 lemma square flop 1 tensor 1}.
    Using \eqref{eq 5 lemma square flop 1 tensor 1} to simplify \eqref{eq 2 lemma square flop 1 tensor 1}, we get that the expression in \eqref{eq 2 lemma square flop 1 tensor 1} becomes
   %  \begin{align}\notag\sum_{1\leq i \leq k-1}\zeta^{-\frac{(n-k)(n-k+1)}{2}-n(k-1-i)}& (-1)^{n-k} \otimes (-1)^{n}\zeta^{\frac{n(n-1)}{2}-1} \\
   %  \notag&= \sum_{i=1}^{k-1} (-1)^k\zeta^{1+2+\ldots +(n-1)- (1+2+ \ldots n-k)-1}\zeta^{-n(k-1)}\zeta^{ni} 1\otimes 1\\
   % \notag &= (-1)^k \zeta^{n-k+1 + \ldots + (n-1)-1}\zeta^{-n(k-1)}\sum_{i=1}^{k-1} \zeta^{ni} \otimes 1\\
   % \notag &=(-1)^k\zeta^{-(1+2+\ldots + (k-1)) -1}(q^{-2}+q^{-4}+ \ldots q^{-2(k-1)}) 1\otimes 1\\
   %  &=(-1)^k\zeta^{-\frac{k(k-1)}{2}-1}q^{-k}[k-1]_q 1\otimes 1 \label{eq 6 lemma square flop 1 tensor 1}
   %  \end{align}
    \begin{align}\notag\sum_{1\leq i \leq k-1}\zeta^{-\frac{(n-k)(n-k+1)}{2}-n(k-1-i)}& (-1)^{n-k} \otimes (-1)^{n}\zeta^{\frac{n(n-1)}{2}-1} \\
   \notag &= (-1)^k \zeta^{-\frac{k(k-1)}{2}-1}\sum_{i=1}^{k-1} \zeta^{ni} \otimes 1\\
    &=(-1)^k\zeta^{-\frac{k(k-1)}{2}-1}q^{-k}[k-1]_q 1\otimes 1 \label{eq 6 lemma square flop 1 tensor 1}
    \end{align}
    as desired.
\end{proof}


\begin{lemma}\label{lemma square flop 1 tensor x1k}
The following equality holds in $\Bimod{\hh{1,k+1}}{}{R^{\hh{1,k+1}}}\otimes_{\hh{1}}\Bimod{}{\hh{0,1}}{R^{\hh{0,1}}}$, for $2\leq k \leq n-2$.
    \begin{align*}
    &\ddm{1,2, k+1}{1,k+1}\left(\ddm{1,2,k+1}{2,k+1}(\Del{1,2}{2}{1})\Del{1,2}{1}{1}\right) \otimes \ddm{0,1,2}{0,1}\left(\ddm{0,1,2}{0,2}(\Del{1,2}{2}{2}x_1^k)\Del{1,2}{1}{2} \right) \\
    =&[k-1]_q(-1)^k\zeta^{\frac{-(k-1)k}{2} -1}q^{-k}(1\otimes x_1^k)
    \;+\;\zeta^{k-1}\ddm{0,1,k+1}{1,k+1}\left( \ddm{0,1,k+1}{0, k+1}(x_1^k)\Del{0,1}{1}{1}\right) \otimes \Del{0,1}{1}{2} 
    \end{align*}
\end{lemma}
\begin{proof}
    We choose $\Del{1,2}{1}{.}$ and $\Del{1,1}{2}{.}$  as in the proof of \cref{lemma square flop 1 tensor 1}, so that the LHS may be written explicitly as 
      \begin{align}
    \sum_{0\leq i,j \leq n-1} \left( \begin{array}{c}
    \ddm{1,2, k+1}{1,k+1}\left(\ddm{1,2,k+1}{2,k+1}\left(e_{n-1-i}(\zeta^{-1}x_1, \zeta^{-2}x_n, \ldots , \zeta^{-(n-1)}x_3)\right)(-1)^{n-1-j}h_{j}(x_2)\right) \vspace{5 pt}\\ \otimes \ddm{0,1,2}{0,1}\left(\ddm{0,1,2}{0,2}((-1)^i\zeta^{n(n-1)/2}h_i(x_2)x_1^k)e_{n-1-j}(\zeta x_3, \zeta^2 x_4, \ldots \zeta^{n-1} x_1) \right) \label{eq 1 lemma square flop 1 tensor x1k}
    \end{array}\right)
    \end{align}
    Let us now simplify the second tensor factor. As in the proof of \cref{lemma square flop 1 tensor 1}, for a given $i,j$, it equals
    \begin{align}
    &(-1)^i\zeta^{n(n-1)/2} \ddm{0,1,2}{0,1}\left( \dd_1(h_i(x_2)x_1^k)e_{n-1-j}(\zeta x_3, \zeta^2 x_4, \ldots , \zeta^{n-1}x_1) \right)\label{eq 2 lemma square flop 1 tensor x1k}
    \end{align}
    \begin{enumerate}
        \item[Case 1:]  Suppose $i=k$. Then $\dd_1(h_i(x_2)x_1^k)=0$, so that there is no contribution to \eqref{eq 1 lemma square flop 1 tensor x1k} from this term.
    \item[Case 2:] Suppose $i>k$. We will not end up needing this case (the other tensor factor vanishes) but we record the simplification without proof for future reference. All the terms with $j \ne i$ vanish, and the term with $j=i$ is equal to $(-1)^n \zeta^{\frac{n(n-1)}{2} - 1} x_1^k$.
    
    
%     Then we have 
%     \begin{align}  \label{eq 3 lemma square flop 1 tensor x1k}
%     (-1)^i&\zeta^{n(n-1)/2} \ddm{0,1,2}{0,1}\left( \dd_1(h_i(x_2)x_1^k)e_{n-1-j}(\zeta x_3, \zeta^2 x_4, \ldots , \zeta^{n-1}x_1) \right)\\
%    \notag &=(-1)^i\zeta^{n(n-1)/2} \ddm{0,1,2}{0,1}\left( (x_1x_2)^k\dd_1(h_{i-k}(x_2))e_{n-1-j}(\zeta x_3, \zeta^2 x_4, \ldots , \zeta^{n-1}x_1) \right)\\
%    \notag &=(-1)^i\zeta^{n(n-1)/2} \ddm{0,1,2}{0,1}\left((x_1x_2)^k\;(-\zeta^{-1}) h_{i-k-1}(x_2, \zeta^{-1}x_1)\; e_{n-1-j}(\zeta x_3, \zeta^2 x_4, \ldots , \zeta^{n-1}x_1) \right)\\ \notag
%     &=(-1)^{i+1}\zeta^{n(n-1)/2-1}x_1^k \sum_{l=0}^{i-k-1} (\zeta^{-1}x_1)^{i-k-1-l} \ddm{0,1,2}{0,1}\left( x_2^{k+l}\; e_{n-1-j}(\zeta x_3, \zeta^2 x_4, \ldots , \zeta^{n-1}x_1) \right)
%     \end{align}
% In the last line we expanded $h_{i-k-1}$ and pulled factors of $x_1$ outside of the Demazure operator.  Note that for the terms above with $i-k-1-l>0$, we'll have no contribution to \eqref{eq 1 lemma square flop 1 tensor x1k} for degree reasons: the degree of $x_1^{k+(i-k-1-l)}$ is $>2k$, forcing the degree of the corresponding first tensor factor to be $<0$ so that the first tensor factor is 0.
%     Hence, the expression in \eqref{eq 3 lemma square flop 1 tensor x1k} simplifies to
%     \begin{align*}
%         &(-1)^{i+1}\zeta^{\frac{n(n-1)}{2}-1}x_1^k \ddm{0,1,2}{0,1}\left( x_2^{i-1}\; e_{n-1-j}(\zeta x_3, \zeta^2 x_4, \ldots , \zeta^{n-1}x_1) \right)\\
%         &=(-1)^{i+1}\zeta^{\frac{n(n-1)}{2}-1}x_1^k \ddm{0,1,2}{0,1}\left( h_{i-1}(x_2)\; \left(\begin{array}{c}
%     e_{n-1-j}(\zeta x_3, \zeta^2 x_4, \ldots , \zeta^{n-2}x_0) \vspace{5pt}\\
%     + \zeta^{n-1}x_1e_{n-2-j}(\zeta x_3, \zeta^2 x_4, \ldots , \zeta^{n-2}x_0)
%     \end{array}\right) \right)
%     \end{align*} 
%     The second term above (one containing $\zeta^{n-1}x_1e_{n-2-j}(\ldots)$) doesn't contribute to \eqref{eq 1 lemma square flop 1 tensor x1k} for degree reasons as before.
%     Moreover, since $k\leq i-1\leq n-2$, we may use \cref{lemma dual basis left ext} to simplify the above expression to
%     \begin{align*}
%         (-1)^{i+1}\zeta^{\frac{n(n-1)}{2}-1}x_1^k  (-1)^{n-2-(i-1)} \delta_{i-1, n-2-(n-1-j)}=(-1)^{n}\zeta^{\frac{n(n-1)}{2}-1}x_1^k \delta_{i-1,j-1}
%     \end{align*}
%     So when $i>k$, we need only consider the terms with $j=i$ in \cref{eq 1 lemma square flop 1 tensor x1k}, with the right tensor factor being $(-1)^n\zeta^{\frac{n(n-1)}{2}-1}x_1^k$.
    \item[Case 3:]  Suppose $0\leq i< k$. Then \eqref{eq 2 lemma square flop 1 tensor x1k} simplifies to 
    \begin{align*}
        (-&1)^i\zeta^{n(n-1)/2} \ddm{0,1,2}{0,1}\left( \dd_1(h_i(x_2)x_1^k)e_{n-1-j}(\zeta x_3, \zeta^2 x_4, \ldots , \zeta^{n-1}x_1) \right)\\
        &=(-1)^i\zeta^{n(n-1)/2} \ddm{0,1,2}{0,1}\left((x_1x_2)^i \dd_1(x_1^{k-i})e_{n-1-j}(\zeta x_3, \zeta^2 x_4, \ldots , \zeta^{n-1}x_1) \right)\\
        &= (-1)^i\zeta^{n(n-1)/2}\ddm{0,1,2}{0,1}\left((x_1x_2)^i h_{k-i-1}(x_1, \zeta x_2)e_{n-1-j}(\zeta x_3, \zeta^2 x_4, \ldots , \zeta^{n-1}x_1) \right)\\
        &=(-1)^i\zeta^{n(n-1)/2}\ddm{0,1,2}{0,1}\left(\sum_{l=0}^{k-1-i}x_1^{i+(k-1-i-l)} \zeta^l x_2^{i+l} e_{n-1-j}(\zeta x_3, \zeta^2 x_4, \ldots , \zeta^{n-1}x_1) \right)\\
        &=(-1)^i\zeta^{n(n-1)/2}\sum_{l=0}^{k-1-i}x_1^{i+(k-1-i-l)} \ddm{0,1,2}{0,1}\left(\zeta^l x_2^{i+l} e_{n-1-j}(\zeta x_3, \zeta^2 x_4, \ldots , \zeta^{n-1}x_1) \right)\\
        &=(-1)^i\zeta^{n(n-1)/2}\sum_{l=0}^{k-1-i}x_1^{k-1-l} \ddm{0,1,2}{0,1}\left(\zeta^l x_2^{i+l}
        \left(\begin{array}{c}
    e_{n-1-j}(\zeta x_3, \zeta^2 x_4, \ldots , \zeta^{n-2}x_0) \vspace{5pt}\\
    + \zeta^{n-1}x_1e_{n-2-j}(\zeta x_3, \zeta^2 x_4, \ldots , \zeta^{n-2}x_0)
    \end{array}\right) \right)
    \end{align*}
    But $i+l\leq i+(k-1-i) =k-1 \leq n-2$, so we can use \cref{lemma dual basis left ext} to simplify the above sum to
    \begin{align}
     \notag (-1)^i\zeta^{n(n-1)/2}\sum_{l=0}^{k-1-i}x_1^{k-1-l}\zeta^l (-1)^{n-2-(i+l)}\left(\delta_{i+l, n-2-(n-1-j)}+\zeta^{n-1}x_1 \delta_{i+l, n-2-(n-2-j)} \right)\\
     \notag =\begin{cases}
         \begin{array}{cl}
           0   & \text{if $j<i$, or $j>k$} \\
           (-1)^i\zeta^{n(n-1)/2}x_1^{i}\zeta^{k-1-i}(-1)^{n-k+1}    & \text{if $j=k$}\\
           (-1)^i\zeta^{n(n-1)/2}x_1^{k+i-j}\zeta^{j-1-i}(-1)^{n-j+1}(1-\zeta^{n})    &\text{if $i<j<k$}\\
           (-1)^i\zeta^{n(n-1)/2}x_1^{k}(-1)^{n-2-j}\zeta^{n-1}    &\text{if $i=j$}
         \end{array}.
     \end{cases}\\
     =\begin{cases}
         \begin{array}{cl}
           0   & \text{if $j<i$, or $j>k$} \\
           (-1)^{n-1+(k-i)}\zeta^{\frac{n(n-1)}{2}-1+(k-i)}x_1^{i}   & \text{if $j=k$}\\
           (-1)^{n-1+(j-i)}\zeta^{\frac{n(n-1)}{2}-1+(j-i)}x_1^{k-(j-i)}(1-q^{-2})    &\text{if $i<j<k$}\\
           (-1)^{n}\zeta^{\frac{n(n-1)}{2}-1}q^{-2}x_1^k    &\text{if $i=j$}
         \end{array}.\label{eq 3.5 lemma square flop 1 tensor x1k}
     \end{cases}
    \end{align}  
    \end{enumerate}

    Now we look at the first tensor factor (for a given $i,j$):
    \begin{align}
        \ddm{1,2, k+1}{1,k+1}\left(\ddm{1,2,k+1}{2,k+1}\left(e_{n-1-i}(\zeta^{-1}x_1, \zeta^{-2}x_n, \ldots , \zeta^{-(n-1)}x_3)\right)(-1)^{n-1-j}h_{j}(x_2)\right) \label{eq 4 lemma square flop 1 tensor x1k}
    \end{align}
     From \eqref{eq left tensor factor square flop} in the proof of \cref{lemma square flop 1 tensor 1}, we have that
     \begin{align*}
         &\dd_{k+2}\ldots \dd_n \dd_1 (e_{n-1-i}(\zeta^{-1}x_1, \zeta^{-2}x_n, \ldots \zeta^{-(n-1)}x_3))\\
         &= \zeta^{-\frac{(n-k)(n-k+1)}{2}}e_{k-1-i}(\zeta^{-(n-k+1)}x_{k+1}, \ldots, \zeta^{-(n-2)}x_4, \zeta^{-(n-1)}x_3),
     \end{align*}
     so that the expression in \eqref{eq 4 lemma square flop 1 tensor x1k} can be written as
     \begin{align}
       \ddm{1,2, k+1}{1,k+1}\left(\zeta^{-\frac{(n-k)(n-k+1)}{2}}e_{k-1-i}(\zeta^{-(n-k+1)}x_{k+1}, \ldots, \zeta^{-(n-2)}x_4, \zeta^{-(n-1)}x_3)(-1)^{n-1-j}h_{j}(x_2)\right) \label{eq 5 lemma square flop 1 tensor x1k}  .
     \end{align}
     Hence we need only consider the case when $0\leq i\leq k-1$,  for which the second tensor factor is explicitly given in Case 3 above (eq \eqref{eq 3.5 lemma square flop 1 tensor x1k}), so that we may further assume that $i\leq j \leq k$.\\
     Since $\ddm{1,2,k+1}{1,k+1}=\dd^{\set{3,4, \ldots, k}}_{\set{2, 3, \ldots, k}}$, by \cref{lemma dual basis left ext}, we have that for $j\leq k-1$, 
     \eqref{eq 5 lemma square flop 1 tensor x1k} is $0$ unless $j=i$.
     Hence the computation from \eqref{eq 5 lemma square flop 1 tensor 1} in the proof of \cref{lemma square flop 1 tensor 1} tells us that the first tensor factor is 
     \begin{align}
    \begin{array}{cc}
    \notag \zeta^{-\frac{(n-k)(n-k+1)}{2}-n(k-1-i)} (-1)^{n-k} \delta_{i,j}  &\text{if $0\leq i,j \leq k-1$,}
    \end{array}
     \end{align}
     and the corresponding term in \eqref{eq 1 lemma square flop 1 tensor x1k} is 
     \begin{align}
         \zeta^{-\frac{(n-k)(n-k+1)}{2}-n(k-1-i)} (-1)^{n-k}\delta_{i,j} \otimes (-1)^{n}\zeta^{\frac{n(n-1)}{2}-1}q^{-2}x_1^k \\
         =(-1)^{k}q^{-2} \zeta^{\frac{n(n-1)}{2}-\frac{(n-k)(n-k+1)}{2}-n(k-1-i)-1}\delta_{i,j} \;\;\;1\otimes x_1^k
         \label{eq 6 lemma square flop 1 tensor x1k} &&\text{if $0\leq i,j \leq k-1$.}
     \end{align}
     
    Now we look at the case when $j=k$, in which case \eqref{eq 5 lemma square flop 1 tensor x1k} can be simplified to
    \begin{align*}
       &\ddm{1,2, k+1}{1,k+1}\left(\zeta^{-\frac{(n-k)(n-k+1)}{2}}e_{k-1-i}(\zeta^{-(n-k+1)}x_{k+1}, \ldots, \zeta^{-(n-2)}x_4, \zeta^{-(n-1)}x_3)(-1)^{n-1-j}h_{k}(x_2)\right) \\
       &=(-1)^{n-1-k}\zeta^{-\frac{(n-k)(n-k+1)}{2}}\dd^{\set{3,4, \ldots, k}}_{\set{2, 3, \ldots, k}}\left(\zeta^{-n(k-1-i)}e_{k-1-i}(\zeta^{k-1}x_{k+1}, \ldots, \zeta^{2}x_4, \zeta x_3)x_2 h_{k-1}(x_2)\right)\\
       &=(-1)^{n-1-k}\zeta^{-\frac{(n-k)(n-k+1)}{2}-n(k-1-i)}\dd^{\set{3,4, \ldots, k}}_{\set{2, 3, \ldots, k}}\left(
       \left(\begin{array}{ll}
       e_{k-i}(\zeta^{k-1}x_{k+1}, \ldots, \zeta^{2}x_4, \zeta x_3,x_2) \\
       -e_{k-i}(\zeta^{k-1}x_{k+1}, \ldots, \zeta^{2}x_4, \zeta x_3)
       \end{array}\right)x_2^{k-1}
       \right).
    \end{align*}
     Since $\dd^{\set{3,4, \ldots, k}}_{\set{2, 3, \ldots, k}}(e_{k-i}(\zeta^{k-1}x_{k+1}, \ldots, \zeta^{2}x_4, \zeta x_3)x_2^{k-1})=0$ for $i\leq k-1$ (by \cref{lemma dual basis left ext}) and $e_{k-i}(\zeta^{k-1}x_{k+1}, \ldots, \zeta^{2}x_4, \zeta x_3,x_2) \in R^{\set{2,3, \ldots ,k}}$, we may further simplify the term above to
     \begin{align}
     \notag &(-1)^{n-1-k}\zeta^{-\frac{(n-k)(n-k+1)}{2}-n(k-1-i)}e_{k-i}(\zeta^{k-1}x_{k+1}, \ldots, \zeta^{2}x_4, \zeta x_3,x_2)\dd^{\set{3,4, \ldots, k}}_{\set{2, 3, \ldots, k}}(
       x_2^{k-1})\\
       &=(-1)^{n-1-k}\zeta^{-\frac{(n-k)(n-k+1)}{2}-n(k-1-i)}e_{k-i}(\zeta^{k-1}x_{k+1}, \ldots, \zeta^{2}x_4, \zeta x_3,x_2),  
     \end{align}
     Using the explicit description of the second tensor factor in \eqref{eq 3.5 lemma square flop 1 tensor x1k}, we have that the corresponding term in \eqref{eq 1 lemma square flop 1 tensor x1k} is:
     \begin{align}
      \notag  &(-1)^{n-1-k}\zeta^{-\frac{(n-k)(n-k+1)}{2}-n(k-1-i)}e_{k-i}(\zeta^{k-1}x_{k+1}, \ldots, \zeta x_3,x_2)\otimes (-1)^{n-1+(k-i)}\zeta^{\frac{n(n-1)}{2}-1+(k-i)}x_1^{i} \\
      &=(-1)^{i}\zeta^{\frac{n(n-1)}{2}-\frac{(n-k)(n-k+1)}{2}-n(k-1-i)-1+(k-i)}e_{k-i}(\zeta^{k-1}x_{k+1}, \ldots,  \zeta x_3,x_2)\otimes x_1^{i} \label{eq 7 lemma square flop 1 tensor x1k}.
     \end{align}
     
    In summary, the only non-zero terms satisfy $0\leq i \leq k-1$ and either $j=i$ or $j=k$. From \eqref{eq 6 lemma square flop 1 tensor x1k}, \eqref{eq 7 lemma square flop 1 tensor x1k}, the expression in \eqref{eq 1 lemma square flop 1 tensor x1k} simplifies to
    \begin{align*}
        &\sum_{\substack{0\leq i \leq k-1\\(j=i)}} (-1)^{k}q^{-2} \zeta^{\frac{n(n-1)}{2}-\frac{(n-k)(n-k+1)}{2}-n(k-1-i)-1} \;\;\;1\otimes x_1^k\hspace{50 pt} 
        \\
        +&\sum_{\substack{0\leq i \leq k-1\\(j=k)}} (-1)^{i}\zeta^{\frac{n(n-1)}{2}-\frac{(n-k)(n-k+1)}{2}-n(k-1-i)-1+(k-i)}e_{k-i}(x_2, \zeta x_3, \ldots, \zeta^{k-1}x_{k+1})\otimes x_1^{i}
    \end{align*}
    The expression above simplifies to
    \begin{align}
       \notag  &\;\;\;\sum_{i=0}^{k-1} (-1)^kq^{-2}\zeta^{-\frac{k(k-1)}{2}-1}\zeta^{ni} 1\otimes x_1^k\\
     \notag   &+\sum_{0\leq i \leq k-1} (-1)^{i}\zeta^{-\frac{k(k-1)}{2}-1+k-i}
        \zeta^{ni}e_{k-i}(x_2, \zeta x_3, \ldots, \zeta^{k-1}x_{k+1})\otimes x_1^{i}\\
    % \notag    =&\;\;\; (-1)^kq^{-2}\zeta^{-\frac{k(k-1)}{2}-1}(1+q^{-2}+\ldots + q^{-2(k-1)}) \;1\otimes x_1^k\\
    %  \notag   &+\sum_{0\leq i \leq k-1} (-1)^{i}\zeta^{-\frac{k(k-1)}{2}-1+k-i}
    %     q^{-2i} e_{k-i}(x_2, \zeta x_3, \ldots, \zeta^{k-1}x_{k+1})\otimes x_1^{i}\\
    % \notag    =&\;\;\; (-1)^k\zeta^{-\frac{k(k-1)}{2}-1}\left((q^{-2}+\ldots + q^{-2(k-1)}) \;1\otimes x_1^k\; + \;q^{-2k}\; 1\otimes x_1^k\right)\\
    %  \notag   &+\sum_{0\leq i \leq k-1} (-1)^{i}\zeta^{-\frac{k(k-1)}{2}-1+k-i}
    %     q^{-2i} e_{k-i}(x_2, \zeta x_3, \ldots, \zeta^{k-1}x_{k+1})\otimes x_1^{i}\\
    =&\;\;\;\notag (-1)^k\zeta^{-\frac{k(k-1)}{2}-1}q^{-k}[k-1]_q\; 1\otimes x_1^k\\
    &+\sum_{0\leq i \leq k} (-1)^{i}\zeta^{-\frac{k(k-1)}{2}-1+k-i}
        q^{-2i} e_{k-i}(x_2, \zeta x_3, \ldots, \zeta^{k-1}x_{k+1})\otimes x_1^{i}. \label{eq 8 lemma square flop 1 tensor x1k}   
    \end{align}
    Comparing the expression above (which is equal the one in \eqref{eq 1 lemma square flop 1 tensor x1k}, which was equal to the LHS of the equation in \cref{lemma square flop 1 tensor x1k}) with the RHS of the equation in \cref{lemma square flop 1 tensor x1k},
    we see that it suffices to show 
    \begin{align}
       \notag \zeta^{k-1}\ddm{0,1,k+1}{1,k+1}&\left( \ddm{0,1,k+1}{0, k+1}(x_1^k)\Del{0,1}{1}{1}\right) \otimes \Del{0,1}{1}{2}\\
        &=\sum_{0\leq i \leq k} (-1)^{i}\zeta^{-\frac{k(k-1)}{2}-1+k-i}
        q^{-2i} e_{k-i}(x_2, \zeta x_3, \ldots, \zeta^{k-1}x_{k+1})\otimes x_1^{i}.\label{eq 9 lemma square flop 1 tensor x1k}
    \end{align}
    Let us make the LHS of the equation above explicit by choosing $
    \Del{0,1}{1}{2}, \Del{0,1,}{1}{1}$ to run over $\set{(-1)^ih_i(x_1)\zeta^{\frac{n(n-1)}{2}}}$, $\set{e_{n-1-i}(\zeta^{-1}x_n, \ldots \zeta^{-(n-1)}x_2)}$ respectively for $0 \leq i \leq n-1$ (using \cref{corollary dual basis right ext} for explicit dual basis). 
    Moreover, 
    \[\ddm{0,1,k+1}{1,k+1}=\dd^{\set{2,3, \ldots,k}\sqcup \set{k+2, \ldots, n-1}}_{\set{2,3, \ldots, k}\sqcup\set{k+2, \ldots, n-1,n}}=\dd^{\set{k+2, \ldots,n-1}}_{\set{k+2,\ldots,n}} ,\]
    \[\ddm{0,1,k+1}{0, k+1}=\dd^{\set{2,3,\ldots,k}\sqcup\set{k+2,\ldots, n-1}}_{\set{1,2,\ldots, k}\sqcup \set{k+2,\ldots,n-1}}=\dd^{\set{2,3,\ldots,k}}_{\set{1,2,\ldots, k}}.\]
    Hence, we have (from \cref{lemma dual basis left ext}) that
    \begin{align*}
        \ddm{0,1,k+1}{0,k+1}(x_1^k)=\dd^{\set{2,3,\ldots,k}}_{\set{1,2,\ldots, k}}(x_1^k)=1.
    \end{align*}
    Hence  the LHS of \eqref{eq 9 lemma square flop 1 tensor x1k} simplifies to
    \begin{align}
        \notag &\zeta^{k-1}\ddm{0,1,k+1}{1,k+1}\left(\Del{0,1}{1}{1}\right) \otimes \Del{0,1}{1}{2}\\
        &=\sum_{0\leq i\leq n-1}\zeta^{k-1}\dd^{\set{k+2, \ldots,n-1}}_{\set{k+2,\ldots,n}}\left( e_{n-1-i}(\zeta^{-1}x_n, \ldots ,\zeta^{-(n-1)}x_2) \right) \otimes (-1)^ih_i(x_1)\zeta^{\frac{n(n-1)}{2}} \label{eq 10 lemma square flop 1 tensor x1k}
    \end{align}
    Since
    \begin{align*}
        &\;\;\;\;\;\;e_{n-1-i}(\zeta^{-1}x_n, \ldots ,\zeta^{-(n-1)}x_2)\\
        &=\sum_{l=0}^{n-1-i}e_l(\zeta^{-1}x_n, \zeta^{-2}x_{n-1}, \ldots ,\zeta^{-(n-k+1)}x_{k+2})e_{n-1-i-l}(\zeta^{-(n-k)}x_{k+1}, \ldots ,\zeta^{-(n-2)}x_3,\zeta^{-(n-1)}x_2),
    \end{align*}
    we have that
    \begin{align*}
      &\dd^{\set{k+2, \ldots,n-1}}_{\set{k+2,\ldots,n}}\left( e_{n-1-i}(\zeta^{-1}x_n, \ldots ,\zeta^{-(n-1)}x_2) \right)  \\
      % &=\dd^{\set{k+2, \ldots,n-1}}_{\set{k+2,\ldots,n}}\left(\sum_{l=0}^{n-1-i}e_l(\zeta^{-1}x_n, \ldots ,\zeta^{-(n-k+1)}x_{k+2})e_{n-1-i-l}(\zeta^{-(n-k)}x_{k+1}, \ldots ,\zeta^{-(n-1)}x_2)  \right)\\
      &=\sum_{l=0}^{n-1-i}e_{n-1-i-l}(\zeta^{-(n-k)}x_{k+1}, \ldots ,\zeta^{-(n-1)}x_2)\dd^{\set{k+2, \ldots,n-1}}_{\set{k+2,\ldots,n}}\left( e_l(\zeta^{-1}x_n, \ldots ,\zeta^{-(n-k+1)}x_{k+2})\right)\\
      &=\sum_{l=0}^{n-1-i}e_{n-1-i-l}(\zeta^{-(n-k)}x_{k+1}, \ldots ,\zeta^{-(n-1)}x_2)\zeta^{-\frac{(n-k)(n-k-1)}{2}}\delta_{l, n-k-1} \hspace{20pt}\text{(Using \cref{corollary dual basis right ext})}\\
      &=\zeta^{-\frac{(n-k)(n-k-1)}{2}}e_{k-i}(\zeta^{-(n-k)}x_{k+1}, \ldots ,\zeta^{-(n-1)}x_2).
    \end{align*}
    Hence the RHS of \eqref{eq 10 lemma square flop 1 tensor x1k} can be simplified to
    \begin{align*}
      &\sum_{0\leq i\leq k}\zeta^{k-1}\zeta^{-\frac{(n-k)(n-k-1)}{2}}e_{k-i}(\zeta^{-(n-k)}x_{k+1}, \ldots ,\zeta^{-(n-1)}x_2) \otimes (-1)^ih_i(x_1)\zeta^{\frac{n(n-1)}{2}}  \\
      &=\sum_{0 \leq i \leq k}(-1)^i\zeta^{\frac{n(n-1)}{2}-\frac{(n-k)(n-k-1)}{2}+k-1}\zeta^{-(n-1)(k-i)}e_{k-i}(\zeta^{k-1} x_{k+1}, \ldots ,x_2) \otimes x_1^i\\
      % &=\sum_{0\leq i \leq k} (-1)^{i}\zeta^{((n-k)+(n-k+1)+\ldots n-1)+k-1-nk+(k-i)}
      %   \zeta^{ni} e_{k-i}(x_2, \zeta x_3, \ldots, \zeta^{k-1}x_{k+1})\otimes x_1^{i}\\
      % &=\sum_{0\leq i \leq k} (-1)^{i}\zeta^{-((k-1)+ \ldots 1)-1+i}
      %   q^{-2i} e_{k-i}(x_2, \zeta x_3, \ldots, \zeta^{k-1}x_{k+1})\otimes x_1^{i}
     &=\sum_{0\leq i \leq k} (-1)^{i}\zeta^{-\frac{k(k-1)}{2}-1+k-i}
        q^{-2i} e_{k-i}(x_2, \zeta x_3, \ldots, \zeta^{k-1}x_{k+1})\otimes x_1^{i}
    \end{align*}
  which is precisely the RHS of \eqref{eq 9 lemma square flop 1 tensor x1k}.
    \end{proof}
    
\begin{lemma}\label{lemma square flop 1 tensor x1k for k=n-1}
The following equality holds in $\Bimod{\hh{0,1}}{}{R^{\hh{0,1}}}\otimes_{\hh{1}}\Bimod{}{\hh{1,2}}{R^{\hh{1,2}}}$.
    \begin{align*}
    &\ddm{0,1,2}{0,1}\left(\ddm{0,1,2}{0,2}(\Del{1,2}{2}{1})\Del{1,2}{1}{1}\right) \otimes \ddm{0,1,2}{0,1}\left(\ddm{0,1,2}{0,2}(\Del{1,2}{2}{2}x_1^{n-1})\Del{1,2}{1}{2} \right) \\
    =& (-1)^{n-1}\zeta^{\frac{-(n-1)(n-2)}{2}-1}q^{-(n-1)}\left([k-1]_q(1\otimes x_1^{n-1})
    \;+\; (-1)^{n-1}\ddm{0,1}{0}(x_1^{n-1})\Del{0,1}{1}{1} \otimes \Del{0,1}{1}{2} \right)
    \end{align*}
\end{lemma}
\begin{proof}
    The LHS simplifies exactly as in \cref{lemma square flop 1 tensor x1k} for $k=n-1$ to the expression in \cref{eq 8 lemma square flop 1 tensor x1k}, so that
    \begin{align}
    \text{LHS} =&\;\;\;\notag (-1)^{n-1}\zeta^{-\frac{(n-1)(n-2)}{2}-1}q^{-(n-1)}[n-2]_q\; 1\otimes x_1^{n-1}\\
    &+\sum_{0\leq i \leq n-1} (-1)^{i}\zeta^{-\frac{(n-1)(n-2)}{2}+n-2-i}
        q^{-2i} e_{n-1-i}(x_2, \zeta x_3, \ldots, \zeta^{n-2}x_{n})\otimes x_1^{i}.
    \end{align}
    Comparing with the RHS, we see that it suffices to show
    \begin{align}
        \notag &\zeta^{\frac{-(n-1)(n-2)}{2}-1}q^{-(n-1)}\ddm{0,1}{0}(x_1^{n-1})\Del{0,1}{1}{1} \otimes \Del{0,1}{1}{2} \\
        &\quad = \sum_{0\leq i \leq n-1} \left(
        (-1)^{i}\zeta^{-\frac{(n-1)(n-2)}{2}+n-2-i}
        q^{-2i}\; e_{n-1-i}(x_2, \zeta x_3, \ldots, \zeta^{n-2}x_{n})\otimes x_1^{i}\right)
    \end{align}
    which is equivalent to showing
    \begin{align}\label{eq 1 lemma square flop 1 tensor x1k for k=n-1}
       q^{-(n-1)}\ddm{0,1}{0}(x_1^{n-1})\Del{0,1}{1}{1} \otimes \Del{0,1}{1}{2} &= \sum_{0\leq i \leq n-1} \left(
        \begin{array}{l}(-1)^{i}\zeta^{n-1-i}
        q^{-2i}\vspace{5 pt}\\
        \quad e_{n-1-i}(x_2, \zeta x_3, \ldots, \zeta^{n-2}x_{n})\otimes x_1^{i}\end{array}\right)
    \end{align}
    By \cref{lemma dual basis left ext}, we know that $\ddm{0,1}{0}(x_1^{n-1})=1$. 
    From \cref{corollary dual basis right ext}, we know that we may choose
    $\Del{0,1}{1}{1}$ to run over $\set{e_{n-1-i}(\zeta^{-1} x_{n}, \zeta^{-2}x_{n-1},\ldots, \zeta^{-(n-1)}x_2 )}_{0 \leq i \leq n-1}$ and $\Del{0,1}{1}{2}$ to run over $\set{(-1)^i\zeta^{\frac{n(n-1)}{2}}x_1^i}_{0\leq i \leq n-1}$.
    So we may expand the LHS of \cref{eq 1 lemma square flop 1 tensor x1k for k=n-1} as
    \begin{align}
        \nonumber q^{-(n-1)}\sum_{0\leq i \leq n-1}&e_{n-1-i}(\zeta^{-1} x_{n}, \zeta^{-2}x_{n-1},\ldots, \zeta^{-(n-1)}x_2 )\otimes (-1)^i\zeta^{\frac{n(n-1)}{2}}x_1^i\\
       \nonumber &=\sum_{0 \leq i\leq n-1}(-1)^i\zeta^{\frac{n(n-1)}{2}}q^{-(n-1)}e_{n-1-i}(\zeta^{-1} x_{n}, \zeta^{-2}x_{n-1},\ldots, \zeta^{-(n-1)}x_2 )\otimes x_1^i\\
       \notag &=\sum_{0\leq i \leq n-1}(-1)^i\zeta^{n(n-1)}\zeta^{-(n-1)(n-1-i)}e_{n-1-i}(\zeta^{n-2} x_{n}, ,\ldots, \zeta x_3, x_2 ) \otimes x_1^i\\
       \notag  &=\sum_{0 \leq i \leq n-1} (-1)^i\zeta^{(n-1)(1+i)}e_{n-1-i}(x_2, \zeta x_3, \ldots, \zeta^{n-2}x_{n} ) \otimes x_1^i\\
       \notag &=\sum_{0 \leq i \leq n-1} (-1)^i\zeta^{n-1-i}q^{-2i}e_{n-1-i}(x_2, \zeta x_3, \ldots, \zeta^{n-2}x_{n} ) \otimes x_1^i
    \end{align}
    which is precisely the RHS of \cref{eq 1 lemma square flop 1 tensor x1k for k=n-1}.
\end{proof}